%~Mouliné par MaN_auto v.0.27.3 2022-08-08 16:38:20
\documentclass[AHL,Unicode,longabstracts,francais,published]{cedram}

\usepackage{float}
\usepackage{tikz}

\DeclareMathOperator{\Grass}{Grass}
\DeclareMathOperator{\covol}{covol}
\DeclareMathOperator{\GL}{GL}
\DeclareMathOperator{\GO}{O}
\DeclareMathOperator{\Supp}{Supp}
\DeclareMathOperator{\diag}{diag}
\DeclareMathOperator{\rg}{rg}
\DeclareMathOperator{\Stab}{Stab}



\usetikzlibrary{shapes.misc}
\tikzset{cross/.style={cross out, draw, minimum size=2*(#1-\pgflinewidth), inner sep=0pt, outer sep=0pt}}







\newcommand{\abs}[1]{\lvert#1\rvert}
\newcommand{\norm}[1]{\lVert#1\rVert}
\newcommand{\bracket}[1]{\langle#1\rangle}
\newcommand{\bcs}{\backslash}
\newcommand{\se}[2]{\{#1,\dots,#2\}}

\newcommand{\R}{\mathbb{R}}
\newcommand{\C}{\mathbb{C}}
\newcommand{\Z}{\mathbb{Z}}
\newcommand{\N}{\mathbb{N}}
\newcommand{\Q}{\mathbb{Q}}
\newcommand{\PP}{\mathbb{P}}
\newcommand{\QQ}{\overline{\mathbb{Q}}}
\newcommand{\NS}{\mathbb{S}}

\newcommand{\bw}{\mathbf{w}}
\newcommand{\be}{\mathbf{e}}
\newcommand{\bv}{\mathbf{v}}
\newcommand{\bp}{\mathbf{p}}

\newcommand{\cA}{\mathcal{A}}
\newcommand{\cM}{\mathcal{M}}
\newcommand{\cF}{\mathcal{F}}

\newcommand{\bbeta}{\widehat{\beta}}
\newcommand{\ggamma}{\widehat{\gamma}}

\newcommand{\eps}{\varepsilon}

\newcommand{\pgcd}{\mathrm{pgcd}}

\makeatletter
\def\editors#1{%
\def\editor@name{#1}
\if@francais
\def\editor@string{Recommand\'e par les \'editeurs \editor@name.}
\else
\def\editor@string{Recommended by Editors \editor@name.}
\fi}
\makeatother

%%%%%%%%%%%%%%%%%%%%%%%%%%%%%%%%%%%%%%%%%%%%%%%%%%%%%%%%%%

\graphicspath{{./figures/}}

\newcommand*{\mk}{\mkern -1mu}
\newcommand*{\Mk}{\mkern -2mu}
\newcommand*{\mK}{\mkern 1mu}
\newcommand*{\MK}{\mkern 2mu}

\hypersetup{urlcolor=purple, linkcolor=blue, citecolor=red}

\newcommand*{\romanenumi}{\renewcommand*{\theenumi}{\roman{enumi}}}
\newcommand*{\Romanenumi}{\renewcommand*{\theenumi}{\Roman{enumi}}}
\newcommand*{\alphenumi}{\renewcommand*{\theenumi}{\alph{enumi}}}
\newcommand*{\Alphenumi}{\renewcommand*{\theenumi}{\Alph{enumi}}}
\let\oldhat\hat
\renewcommand*{\hat}[1]{\mathchoice{\widehat{#1}}{\widehat{#1}}{\oldhat{#1}}{\oldhat{#1}}}
\let\oldexists\exists
\renewcommand*{\exists}{\mathrel{\oldexists}}
\let\oldforall\forall
\renewcommand*{\forall}{\mathrel{\oldforall}}
%%%%%%%%%%%%%%%%%%%%%%%%%%%%%%%%%%%%%%%%%%%%%%%%%%%%%%%

\title{Approximation diophantienne sur les quadriques}
\alttitle{Diophantine approximation on quadrics}

\subjclass{11J83, 11J87, 37A17}
\keywords{espaces de réseaux, points rationnels, groupes orthogonaux, théorème du sous-espace}


\author[\initial{N.} \lastname{de Saxcé}]{\firstname{Nicolas} \lastname{de Saxcé}}
\address{CNRS -- Université Sorbonne Paris Nord\\
LAGA / UMR 7539\\
99 avenue Jean-Baptiste Clément,\\
93430 Villetaneuse (France)}
\email{desaxce@math.univ-paris13.fr}


\begin{abstract}
Nous étudions l'exposant diophantien d'un point $x$ d'une hypersurface quadratique. Nous montrons notamment un analogue du théorème de Thue--Siegel--Roth, c'est-à-dire une formule pour l'exposant diophantien d'un point algébrique, et un analogue du résultat de Kleinbock et Margulis sur l'extrémalité des sous-variétés non dégénérées de l'espace affine.
\end{abstract}

\begin{altabstract}
We study the Diophantine exponent of a point $x$ on a quadric hypersurface. We show in particular an analogue of the Thue--Siegel--Roth theorem, that is to say a formula for the Diophantine exponent of an algebraic point, and an analogue of the result of Kleinbock and Margulis on the extremality of non-degenerate analytic manifolds in the of affine space.
\end{altabstract}


\datereceived{2021-07-09}
\daterevised{2022-03-30}
\dateaccepted{2022-05-12}

\editors{S. Gou\"ezel et X. Caruso}
%%%%%%%%%%%%%%%%%%%%%%%%%%%%%%%%%%%%%%%%%%%%%%%%%%%%

\dateposted{2022-11-14}
\begin{document}
\maketitle


%\newpage

\section{Introduction}


Suivant une terminologie introduite par Minkowski au début du XX$^{\rm ème}$ siècle, on appelle aujourd'hui \og~approximation diophantienne~\fg{} l'étude des approximations \emph{rationnelles} des points \emph{irrationnels} d'un espace. Traditionnellement, on cherche à approcher un point $x\in\R^n$ par un point rationnel $v=\frac{\bp}{q}\in\Q^n$, et pour évaluer la qualité de l'approximation, on compare la distance $d(x,v)$ au dénominateur $q$ de l'approximation rationnelle.

Lorsque le point $x$ est choisi sur une sous-variété $X$ fixée dans $\R^n$, on parle d'approximation diophantienne sur les sous-variétés. Si de plus on ne s'intéresse qu'aux approximations rationnelles qui appartiennent à $X$, on parle d'approximation diophantienne \emph{intrinsèque}. Cela permet de distinguer ce problème de celui de l'approximation diophantienne par des points de l'espace ambiant $\R^n$, non nécessairement dans $X$. Notons toutefois que ces deux problèmes sont reliés~: par exemple, lorsque la variété $X$ est une hypersurface quadrique rationnelle, il est facile de voir que si $x\in X$, tout point rationnel $\frac{\bp}{q}$ tel que $d(x,\frac{\bp}{q})\leq q^{-2+\eps}$ doit appartenir à $X$, si $q$ est assez grand. L'utilisation de cette observation élémentaire par Dru{\c t}u~\cite{drutu} pour obtenir un analogue du théorème de Jarník sur les quadriques a d'ailleurs certainement participé au regain d'intérêt récent pour ces problèmes d'approximations diophantienne intrinsèque. Ainsi, à l'aide des outils de la dynamique homogène, Kleinbock et Merrill~\cite{kleinbock-merrill} ont pu montrer que plusieurs théorèmes classiques d'approximation diophantienne --- théorèmes de Dirichlet, de Khintchine, etc. --- admettent des analogues naturels pour l'approximation intrinsèque sur la sphère unité dans $\R^n$, et ces résultats ont même été généralisés dans~\cite{fkms_quadric} pour s'appliquer à une hypersurface quadrique rationnelle quelconque.

Nous voulons ici poursuivre leur étude, en abordant en particulier deux cas: celui où le point $x$ dans la quadrique $X$ est à coordonnées dans le corps $\QQ$ des nombres réels algébriques, et celui où $x$ est choisi aléatoirement sur une sous-variété analytique strict de $X$. Cela nous mènera à énoncer et démontrer deux résultats, le premier est un analogue du célèbre théorème de Roth~\cite{roth} sur l'exposant d'approximation des points algébriques par les rationnels, et le second s'inspire des résultats de Kleinbock et Margulis~\cite{kleinbockmargulis} sur l'extrémalité des sous-variétés non dégénérées de $\R^n$.

%\bigskip

Les premiers résultats d'approximation diophantienne sur les quadriques ont souvent été énoncés pour les hypersurfaces quadriques affines. Cependant, pour notre approche, il est mieux adapté de se placer dans un cadre projectif, ce que nous ferons donc dans tout cet article. Les deux cadres sont essentiellement équivalents.

Soit $F$ une forme quadratique rationnelle sur $\R^d$, et $X\subset\PP^{d-1}$ la quadrique associée:
\[
X = \left\{x\in\PP^{d-1}\, \middle|\, F(x)=0\right\}.
\]
Sur l'ensemble $X(\Q)$ des points rationnels de $X$, on dispose de la hauteur induite par la hauteur usuelle sur $\PP^{d-1}(\Q)$: si $v\in\PP^{d-1}(\Q)$ a pour coordonnées homogènes $[v_1:\dots:v_d]$, avec $v_i\in\Z$ et $\pgcd(v_1,\,\dots,\,v_d)=1$, alors
\[
H(v) = \max_{1\,\leq\,i\,\leq\,d} \left|v_i\right|.
\]
Nous munissons aussi $X(\R)$ de la distance induite par $\PP^{d-1}(\R)$:
\[
d(x,y) = \frac{\|x\wedge y\|}{\|x\|\,\|y\|}.
\]
Dans cette formule, la norme considérée sur $\R^d$ est celle de la structure euclidienne canonique, et dans le membre de droite, nous avons noté abusivement $x$ et $y$ pour des représentants dans $\R^d$ de $x$ et $y$. Nous supposons en général que $X(\Q)$ est dense dans $X(\R)$, et définissons l'\emph{exposant diophantien} d'un point $x\in X(\R)$ par

\[
\beta(x) = \inf\left\{\beta>0\, \middle|\, \exists c>0:\,\forall v\in X(\Q),\ d(x,v)\geq cH(v)^{-\beta}\right\}.
\]
Si l'ensemble de droite est vide, ce qui est le cas par exemple lorsque $x\in X(\Q)$, on pose $\beta(x)=+\infty$. Le théorème suivant, dû à Fishman, Kleinbock, Merrill et Simmons~\cite{fkms_quadric} donne la valeur minimale de l'exposant $\beta(x)$ pour $x\in X(\R)$, et montre que cette valeur est atteinte pour presque tout point dans $X(\R)$ pour la mesure de Lebesgue.

\begin{theo}[Fishman--Kleinbock--Merrill--Simmons]
Soit $F$ une forme quadratique rationnelle sur $\R^d$, et $X$ la quadrique associée. On suppose que $X(\Q)\setminus\ker F\neq\varnothing$. Alors,
\begin{enumerate}
\item pour tout $x\in X(\R)$, $\beta(x)\geq 1$;
\item pour presque tout $x\in X(\R)$, $\beta(x) = 1$.
\end{enumerate}
\end{theo}

Le but de cet article est d'étudier l'exposant $\beta(x)$ dans les deux cas suivants~:
\begin{itemize}
\item $x$ est un point algébrique de $X$, i.e. admet des coordonnées homogènes dans un corps de nombres;
\item $x$ est choisi aléatoirement sur une sous-variété analytique $\cM$ de $X$.
\end{itemize}
Notre premier théorème est un analogue pour les quadriques du théorème de Thue--Siegel--Roth, qui donne en particulier une condition nécessaire et suffisante pour que du point de vue de l'exposant diophantien, un point algébrique $x$ se comporte comme un point choisi aléatoirement pour la mesure de Lebesgue. Rappelons qu'un sous-espace $W$ est dit \emph{totalement isotrope} pour une forme quadratique $F$ sur $\R^d$, si la restriction de $F$ à $W$ est nulle. Si $X=\{F=0\}\subset\PP^{d-1}$ est la quadrique associée à $F$, les sous-espaces projectifs totalement isotropes de $X$ sont les sous-espaces projectifs inclus dans $X$. Ci-dessous et dans toute la suite, on note $\QQ\subset\R$ le corps des nombres algébriques réels.

\begin{theo}[Exposant d'un point algébrique]\label{roth_quadrique}
Soit $X$ une quadrique rationnelle dans $\PP^{d-1}$. Si $x\in X(\QQ)$ n'est inclus dans aucun sous-espace projectif rationnel totalement isotrope, alors $\beta(x)=1$. Plus généralement, pour tout $x$ dans $X(\QQ)$,
\[
\beta(x) = 1 + \frac{1}{k_x},
\]
où $k_x$ est la dimension minimale d'un sous-espace projectif rationnel totalement isotrope contenant $x$. (Par convention, on pose $k_x=+\infty$ s'il n'existe pas de tel sous-espace projectif contenant $x$.)
\end{theo}

Lorsque $X$ est une sphère, un sous-espace totalement isotrope est toujours réduit à un point, donc $k_x=+\infty$ pour tout $x$ irrationnel. La formule pour l'exposant est donc particulièrement simple dans ce cas.

\begin{coro}[Théorème de Roth pour la sphère]
Pour tout point $x$ irrationnel dans $\NS^n(\QQ)$, $\beta(x)=1$.
\end{coro}

Nous nous tournons maintenant vers la situation où le point $x$ est choisi aléatoirement sur une sous-variété analytique $\cM$ dans $X(\R)$. La variété $\cM$ sera toujours munie de la restriction à $\cM$ de la mesure de Hausdorff en dimension $m=\dim\cM$, que nous appellerons parfois \og~mesure de Lebesgue sur $\cM$~\fg{}. Le théorème ci-dessous découle d'une observation plus générale de Kleinbock~\cite{kleinbock_dichotomy} sur le taux de fuite des orbites diagonales dans l'espace des réseaux. Nous en donnerons néanmoins la démonstration à la partie~\ref{sec:existence_exposant}.

\begin{theo}[Exposant d'une variété analytique]\label{exexpi}
Soit $X$ une quadrique rationnelle dans $\PP^{d-1}$, et $\cM$ une sous-variété analytique connexe de $X(\R)$. On note $\bbeta(\cM)=\inf_{x\,\in\,\cM}\beta(x)$. Alors, pour presque tout $x$ dans $\cM$,
\[
\beta(x) = \bbeta(\cM).
\]
\end{theo}

On remarquera que cette observation de Kleinbock, combinée avec le théorème ci-dessus pour les points algébriques, donne déjà une formule pour l'exposant d'un point choisi aléatoirement sur une sous-variété analytique $\cM$ dont l'adhérence de Zariski est définie sur $\QQ$.

\begin{theo}[Exposant d'un ensemble algébrique défini sur $\QQ$]\label{algi}
Soit $X$ une quadrique rationnelle dans $\PP^{d-1}$ et $\cM$ une sous-variété analytique connexe de $X(\R)$. Si $\cM$ n'est incluse dans aucun sous-espace projectif rationnel totalement isotrope, alors pour presque tout $x\in\cM$, $\beta(x)=1$. Plus généralement, si l'adhérence de Zariski de $\cM$ est définie sur $\QQ$, alors, pour presque tout $x\in\cM$,
\[
\beta(x) = 1 + \frac{1}{k_{\cM}},
\]
où $k_{\cM}$ est la dimension minimale d'un sous-espace projectif rationnel totalement isotrope contenant $\cM$.
\end{theo}


Le second objectif de cet article est d'obtenir une formule pour l'exposant presque sûr $\bbeta(\cM)$ d'une variété analytique qui n'est pas nécessairement définie sur $\QQ$. Bien sûr, cette formule met en jeu certaines propriétés diophantiennes de $\cM$. Pour l'énoncer, nous aurons besoin d'une hauteur sur l'ensemble des sous-espaces rationnels projectifs de $\PP^{d-1}$. Celle-ci généralise la hauteur usuelle sur $\PP^{d-1}(\Q)$ et se définit à partir du plongement de Pl\"ucker~; nous renvoyons à la partie~\ref{sec:formule_exposant} pour plus de précisions. Étant donnée une partie $S$ de $X(\R)$, et $i\in\se{1}{d}$, nous pouvons alors définir le $i$-ème exposant diophantien de $S$ de la façon suivante~:
\[
\beta_i(S) = \inf\left\{\beta>0\ \middle|
\begin{split}
\exists c>0:\,\forall W\ \mbox{totalement isotrope rationnel}\\
\mbox{de dimension}\ i,\ d(S,W)\geq cH(W)^{-\beta}
\end{split}
\right\}.
\]
Dans l'égalité ci-dessus, $d(S,W)=\sup_{x\,\in\,S} d(x,W)$.

\begin{rema}
Si $S=\{x\}$, l'exposant $\beta_0(x)$ n'est autre que l'exposant diophantien $\beta(x)$ défini au début de cette introduction. Par ailleurs, si $S$ n'est inclus dans aucun sous-espace totalement isotrope de dimension~$i$, alors $\beta_i(S)=0$. Plus généralement, on se convainc sans peine que les exposants $\beta_i(S)$ ne dépendent que du plus petit sous-espace totalement isotrope $W_S$ contenant $S$; \emph{cf.} lemme~\ref{her}.
\end{rema}

\begin{theo}[Hérédité et formule pour l'exposant]\label{formexpi}
Soit $X$ une quadrique rationnelle dans $\PP^{d-1}$, et $\cM$ une sous-variété analytique connexe de $X(\R)$. Si $\cM$ n'est incluse dans aucun sous-espace projectif totalement isotrope, alors
\[
\bbeta(\cM) = 1.
\]
Plus généralement, si $W_{\cM}$ désigne le plus petit sous-espace totalement isotrope contenant $\cM$, alors l'exposant $\bbeta(\cM)$ est donné par la formule
\[
\bbeta(\cM) = \max\left(1,\max_{0\,\leq\,i\,\leq\, d-1} \frac{(i+1)\beta_i(W_{\cM})}{1+i\beta_i(W_{\cM})}\right).
\]
En particulier, $\bbeta(\cM)$ est entièrement déterminé par le sous-espace $W_{\cM}$.
\end{theo}

\begin{exem}
Dans le cas où $\cM=\{x\}$ est réduite à un singleton, on a par définition $\bbeta(\{x\})=\beta_0(x)$. Nous verrons dans la démonstration de la proposition~\ref{dirichlet-fort} ci-dessous que pour chaque $i\geq 1$, $\beta_0(x) \geq \frac{(i+1)\beta_i(x)}{1+i\beta_i(x)}$. Le théorème est donc trivial dans ce cas.
\end{exem}

\begin{exem}
Si $\cM$ n'est pas réduite à un singleton, le maximum dans la formule pour $\bbeta(\cM)$ n'est pas nécessairement atteint pour $i=\dim W_{\cM}$. Par exemple, si $\cM$ est une droite projective prise aléatoirement dans un plan rationnel totalement isotrope fixé, on vérifie aisément que $\beta_0(\cM)=0$, $\beta_1(\cM)=3/2$, et $\beta_2(\cM)=\infty$. Donc le maximum est atteint pour $i=2$, et $\bbeta(\cM)=\frac{3}{2}$. En revanche, si $\cM$ est une droite projective rationnelle, alors $\beta_0(\cM)=0$ et $\beta_1(\cM)=\infty$, donc le maximum est atteint pour $i=1$, et $\bbeta(\cM)=2$.
\end{exem}

Comme corollaire de ce théorème, on peut affaiblir l'hypothèse du théorème~\ref{algi} sur les sous-variétés définies sur $\QQ$: il suffit de supposer que le plus petit sous-espace totalement isotrope contenant $\cM$ est défini sur $\QQ$.

\begin{coro}\label{petitplus}
Soit $X$ une quadrique rationnelle dans $\PP^{d-1}$ et $\cM$ une sous-variété analytique connexe de $X(\R)$. On suppose que le plus petit sous-espace totalement isotrope contenant $\cM$ est défini sur $\QQ$. Alors, pour presque tout $x\in\cM$,
\[
\beta(x) = 1 + \frac{1}{k_{\cM}},
\]
où $k_{\cM}$ est la dimension minimale d'un sous-espace projectif rationnel totalement isotrope contenant $\cM$.
\end{coro}

Les résultats énoncés ci-dessus sont tous démontrés à l'aide de la correspondance mise en évidence par Kleinbock, Merrill, Fishman et Simmons~\cite{kleinbock-merrill,fkms_quadric} qui relie l'exposant diophantien d'un point $x\in X(\R)$ au comportement asymptotique d'une orbite diagonale bien choisie dans l'espace des réseaux. Par ailleurs, pour étudier les propriétés diophantiennes des points algébriques, nous utiliserons une interprétation du théorème du sous-espace de Schmidt en termes d'orbites diagonales dans l'espace des réseaux tirée d'un travail en commun avec Emmanuel Breuillard~\cite{breuillardsaxce}, mais dont nous donnerons ici une démonstration alternative. Enfin, le théorème~\ref{formexpi} se démontre en suivant la méthode utilisée par Kleinbock dans~\cite{kleinbock_anextension} dans son étude des exposants diophantiens des sous-espaces, une fois démontré un énoncé de non divergence quantitative adapté au groupe orthogonal d'une forme quadratique, le théorème~\ref{nondiv} ci-dessous. Nous insistons sur le fait que dans cet énoncé, les hypothèses ne concernent que les sous-réseaux totalement isotropes de $\Z^d$, ce qui est crucial pour l'application au théorème~\ref{formexpi}. Ce point important empêche de déduire le résultat des énoncés de non divergence de Kleinbock et Margulis~\cite{kleinbockmargulis,kleinbock_anextension}, et même, leur démonstration ne semble pas s'adapter. Nous démontrons donc ce théorème à la partie~\ref{sec:nd} avec une approche sensiblement différente, qui peut aussi donner une autre démonstration des résultats de Kleinbock et Margulis pour $\GL_d$, exposée en détail dans~\cite{saxce_hnnondivergence}.



\section{Orbites diagonales dans l'espace des réseaux}\label{sec:dani}

Nous exposons dans cette partie la \og~correspondance de Dani~\fg{} pour les quadriques, qui relie l'exposant diophantien $\beta(x)$ d'un point $x$ dans $X(\R)$ au taux de fuite d'une certaine orbite dans l'espace des réseaux de $\R^d$. Cette correspondance a été démontrée par Kleinbock et Merrill~\cite{kleinbock-merrill} dans le cas de la sphère, puis par Fishman, Kleinbock, Merrill et Simmons~\cite{fkms_quadric} pour une quadrique générale.


Comme précédemment, $F$ est une forme quadratique rationnelle sur $\R^d$, et $X$ la quadrique associée dans $\PP^{d-1}$. Pour éviter des cas dégénérés, nous supposerons en général que $X(\Q)\setminus\ker F$ est non vide, de sorte que dans une certaine base rationnelle $(e_i)_{1\leq i\leq d}$ de $\R^d$, la forme quadratique $F$ peut s'écrire
\[
F\left(x_1,\,\dots,\,x_d\right) = 2x_1x_d + G\left(x_2,\dots,x_{d-1}\right),
\]
où $G$ est une forme quadratique rationnelle sur $\R^{d-2}$.

Par le théorème de Witt, si $x$ est un point de $X\setminus\ker F$, il existe un élément $u_x$ dans le groupe orthogonal $O(F)$ tel que $u_x\cdot x = [e_1]$. On considère alors un flot diagonal $(g_t^x)_{t\,>\,0}$ défini dans la base $(e_i)_{1\,\leq\,i\,\leq\,d}$ par
\[
g_t^x = \diag\left(e^{-t},1,\,\dots,\,1,e^t\right)\cdot u_x.
\]
On rappelle qu'un réseau dans $\R^d$ est un sous-groupe discret de rang maximal; et que la \emph{systole} d'un réseau $\Delta$ est par définition
\[
\lambda_1(\Delta) = \min\left\{\|v\|\ ;\ v\in\Delta\setminus\{0\}\right\}.
\]
Le \emph{taux de fuite} $\gamma(x)$ de l'orbite $g_t^x\Z^d$ dans l'espace des réseaux est
\[
\gamma(x) = \inf\left\{\gamma>0\, \middle|\, \exists c>0:\,\forall t>0,\, \lambda_1\left(g_t^x\Z^d\right)\geq ce^{-\gamma t}\right\}.
\]
La proposition ci-dessous, due à Kleinbock et Merrill~\cite[Theorem~1.5]{kleinbock-merrill} dans le cas de la sphère, et à Fishman--Kleinbock--Merrill--Simmons~\cite[Lemma~7.2]{fkms_quadric} dans le cas général, permet d'obtenir l'exposant diophantien $\beta(x)$ à partir du taux de fuite $\gamma(x)$.

\begin{prop}[Correspondance de Dani pour les quadriques]\label{dani}
Soit $F$ une forme quadratique rationnelle sur $\R^d$ et $X$ la quadrique associée dans $\PP^{d-1}$. Pour tout $x$ dans $X(\R)\setminus\ker F$,
\[
\beta(x) = \frac{1}{1-\gamma(x)}.
\]
\end{prop}

\begin{proof}
Nous suivons essentiellement l'argument présenté dans\linebreak\cite[Lemma~7.2]{fkms_quadric}. Dans cette démonstration, les constantes impliquées dans la notation de Vinogradov $\ll$ et $\asymp$ peuvent dépendre de $x$. Pour $i=1,\,\dots,\,d$, soit $e_{i,x}=\linebreak u_x^{-1}e_i$. Si $v$ est un vecteur de $\R^d$, on écrit

\[
v= \sum_{i=1}^d v_x^{(i)}e_{i,x},
\]
de sorte que pour $t>0$, dans la base $(e_{i,x})_{1\,\leq\,i\,\leq\,d}$,
\begin{equation}\label{gtv}
g_t^x v =
\begin{pmatrix}
e^{-t}v_x^{(1)} \\
v_x^{(2)}\\
\vdots\\
v_x^{(d-1)}\\
e^t v_x^{(d)}
\end{pmatrix}
\end{equation}
Supposons maintenant que pour $\beta>0$, on ait $d(x,v)\leq H(v)^{-\beta}$, pour $v\in X(\Q)$ arbitrairement grand. Nous identifions $v$ à l'un de ses représentants primitifs dans $\Z^d$. Alors, 
\[
H(v)=\|v\|, d(x,v)\asymp\frac{\max_{2\,\leq\,i\,\leq\,d}v_x^{(i)}}{\|v\|},
\]
et
\[
0 = F(v) = F(u_xv) = 2v_x^{(1)}v_x^{(d)} + G\left(v_x^{(2)},\,\dots,\,v_x^{(d-1)}\right).
\]
Par conséquent,
\[
\left|v_x^{(d)}\right|=\frac{\left|G\left(v_x^{(2)},\,\dots,\,v_x^{(d-1)}\right)\right|}{2\left|v_x^{(1)}\right|}
\ll \frac{\max\limits_{2\,\leq\,i\,\leq \,d} \left|v_x^{(i)}\right|^2}{\|v\|}
\ll \|v\|^{-2\beta+1}.
\]
Choisissant $t>0$ tel que $e^{t}=\|v\|^{\beta}$, l'égalité~\eqref{gtv} implique $\|g_t^xv\|\ll e^{-(1-\frac{1}{\beta})t}$. Donc $\gamma(x)\geq 1-\frac{1}{\beta(x)}$.

Réciproquement, supposons que pour $\gamma>0$, on a $\|g_t^xv\|\leq e^{-\gamma t}$ pour $t>0$ arbitrairement grand et $v\in\Z^d$. Comme $F(v)=F(g_t^xv)$ est un entier borné en valeur absolue par $e^{-\gamma t}<1$, on doit avoir $F(v)=0$, donc le vecteur $v$ correspond à un point rationnel sur la quadrique $X$. De plus, par~\eqref{gtv}, $|v_x^{(1)}|\leq e^{(1-\gamma)t}$, $|v_x^{(d)}|\leq e^{-(1+\gamma)t}$ et $|v_x^{(i)}|\leq e^{-\gamma t}$ pour $i=2,\,\dots,\,d-1$. Cela montre que $\|v\|\asymp |v_x^{(1)}|$, d'où $\|v\|\ll e^{(1-\gamma)t}$ puis
\[
d(x,v) \asymp \frac{\max\limits_{2\,\leq\,i\,\leq\,d-1}\left|v_x^{(i)}\right|}{\|v\|} \leq \frac{e^{-\gamma t}}{\|v\|} \leq \|v\|^{-\frac{1}{1-\gamma}}.
\]
Donc $\beta(x)\geq\frac{1}{1-\gamma(x)}$.
\end{proof}

La proposition ci-dessus a pour premier corollaire une borne inférieure sur l'exposant diophantien d'un point quelconque de $X(\R)$; c'est l'analogue du principe de Dirichlet pour l'approximation diophantienne dans $\PP^{n}$. Pour un énoncé plus précis, on renvoie le lecteur à~\cite[Theorem~4.1]{kleinbock-merrill} ou~\cite[Theorem~5.1]{fkms_quadric}, d'où ce résultat provient.

\begin{coro}[Principe de Dirichlet pour les quadriques]\label{dirichlet}
Soit $X$ une quadrique rationnelle telle que $X(\Q)$ soit dense dans $X(\R)$. Alors, pour tout $x\in X(\R)$, $\beta(x)\geq 1$.
\end{coro}


\begin{proof}
D'après la proposition~\ref{dani}, il suffit de montrer que $\gamma(x)\geq 0$. Pour cela, on observe que le covolume du réseau $g_t^x\Z^d$ est indépendant de $t$. Le premier théorème de Minkowski permet donc d'en majorer la systole, indépendamment de $t$, ce qui implique en particulier que $\gamma(x)\geq 0$.
\end{proof}



\section{Le théorème du sous-espace fort de Schmidt}\label{sec:algebrique}

À l'aide de la correspondance expliquée à la partie précédente, nous démontrons maintenant la formule pour l'exposant diophantien d'un point algébrique dans une quadrique rationnelle.

%\bigskip

Notre argument repose sur une version forte du théorème du sous-espace de Schmidt~\cite[Theorem~3A, page~163]{schmidt_da}, dont nous rappelons l'énoncé ci-dessous. Étant donné un réseau $\Delta$ dans $\R^d$, les \emph{minima successifs} $\lambda_1(\Delta)\leq \lambda_2(\Delta)\leq\dots\leq\lambda_d(\Delta)$ sont définis par
\[
\lambda_k(\Delta) = \inf\left\{\lambda>0\, \middle|\, \rg(B(0,\lambda)\cap\Delta)\geq k\right\},
\]
où l'on note $\rg(E)$ le rang linéaire d'une famille $E$ d'éléments de $\R^d$.

\begin{rema}
Le premier minimum $\lambda_1(\Delta)$ est égal à la systole $\lambda_1(\Delta)$ que nous avons étudié à la partie précédente. Les notations sont cohérentes.
\end{rema}

Dans la suite, $\R^d$ sera muni de sa structure rationnelle canonique, et nous dirons qu'un réseau $\Delta\leq\R^d$ est \emph{algébrique} s'il admet une base constituée de vecteurs à coordonnées algébriques dans $\R^d$. Les sous-groupes d'un réseau $\Delta$ seront appelés \emph{sous-réseaux} de $\Delta$; ce sont des sous-groupes discrets de $\R^d$, de rang non nécessairement maximal.

\begin{theo}[Théorème du sous-espace fort de Schmidt]
Soit $(a_t)_{t\,>\,0}$ un flot diagonal dans $\GL_d(\R)$ et $\Delta$ un réseau algébrique dans $\R^d$. Supposons qu'il existe une sous-suite $(t_n)_{n\,\geq\,0}$ convergeant vers $+\infty$, un entier $k\in\se{1}{d-1}$ et $\eps>0$ tel que pour chaque $n$, $\lambda_k(a_{t_n}\Delta)\leq e^{-\eps t_n}\lambda_{k+1}(a_{t_n}\Delta)$. Alors il existe une sous-suite $(t_{n_p})_{p\,\geq\,0}$ et un sous-réseau $T\leq\Delta$ de rang $k$ tel que pour tout $p$, les $k$ premiers minima de $a_{t_{n_p}}\Delta$ sont atteints dans $a_{t_{n_p}}T$.
\end{theo}

Le théorème du sous-espace fort permet de décrire le comportement asymptotique des minima successifs des réseaux d'une orbite diagonale $(a_t\Delta)_{t\,>\,0}$, lorsque $\Delta$ est un réseau algébrique. Le théorème suivant peut aussi se déduire directement de l'énoncé usuel du théorème du sous-espace de Schmidt, auquel il est d'ailleurs équivalent; cela est expliqué en détail dans~\cite{breuillardsaxce}.

\begin{theo}\label{drapeau}
Soit $(a_t)_{t\,>\,0}$ un flot diagonal dans $\GL_d(\R)$ et $\Delta$ un réseau algébrique dans $\R^d$. Pour chaque $i\in\se{1}{d}$, la limite $\Lambda_i = \lim_{t\,\to\,\infty} \frac{1}{t}\log\lambda_i(a_t\Delta)$ existe. De plus, si les indices $i_1<\dots<i_s$ sont choisis de sorte que
\[
\Lambda_1=\dots=\Lambda_{i_1}<\Lambda_{i_1+1}=\dots=\Lambda_{i_2}<\dots<\Lambda_{i_s+1}=\dots=\Lambda_d,
\]
alors il existe un drapeau partiel $0=T_0<T_1<\dots<T_{s+1}=\R^d$ de sous-espaces $\Delta$-rationnels tel que pour chaque $\ell\in\se{1}{s}$,
\begin{itemize}
\item $\dim T_\ell = i_\ell$;
\item pour tout $t>0$ assez grand, $a_tT_\ell$ contient les $i_\ell$ premiers minima de $a_t\Delta$.
\end{itemize}
\end{theo}

Pour la démonstration de ce théorème, il sera pratique d'utiliser les puissances extérieures de $\R^d$, notées $\wedge^k\R^d$, pour $k\in\se{1}{d}$. Rappelons que $\wedge^k\R^d$ est un espace vectoriel de dimension $\binom{d}{k}$. Si $I=\{i_1<\dots<i_k\}$ est une partie de $\se{1}{d}$ à $k$ éléments, on note
\[
\be_I=e_{i_1}\wedge\dots\wedge e_{i_k}.
\]
Les vecteurs $\be_I$, avec $I\subset\se{1}{d}$ et $|I|=k$, forment une base de $\wedge^k\R^d$, et l'on munit $\wedge^k\R^d$ de l'unique structure euclidienne telle que $(\be_I)_I$ soit orthonormée.

Si $W$ est un sous-espace vectoriel de $\R^d$, nous dirons qu'un vecteur $\bw$ de $\wedge^k\R^d$ \emph{représente} $W$ s'il existe une base $(v_1,\dots,v_k)$ de $W$ telle que $\bw=v_1\wedge\dots\wedge v_k$. Tous les vecteurs représentant un sous-espace $W$ sont colinéaires, et cela permet de définir le \emph{plongement de Pl\"ucker} de la grassmannienne $\Grass(k,d)$ des $k$-plans de $\R^d$ dans l'espace projectif $\PP(\wedge^k\R^d)$.

Si $\Lambda$ est un sous-groupe discret de $\R^d$, nous dirons qu'un vecteur $\bw$ \emph{représente} $\Lambda$ s'il existe des vecteurs $v_1,\,\dots,\,v_k$ tels que $\Lambda=\Z v_1\oplus\dots\oplus\Z v_k$ et $\bw=v_1\wedge\dots\wedge v_k$. Si $\bw$ représente $\Lambda$, la direction de $\bw$ donne le sous-espace vectoriel $W_\Lambda$ engendré par $\Lambda$, tandis que la norme de $\bw$ est égale au covolume de $\Lambda$ dans $W_\Lambda$~:
\[
\|\bw\| = \covol_{W_\Lambda}(\Lambda).
\]
En particulier, au signe près, il existe un unique vecteur qui représente $\Lambda$. Ce formalisme permet de définir le taux de dilatation d'un sous-espace par un flot diagonal de la façon suivante.

\begin{defi}[Taux de dilatation d'un sous-espace]
Soit $(a_t)_{t\,>\,0}$ un flot diagonal dans $\GL_d(\R)$. Si $W$ est un sous-espace de dimension $k$ dans $\R^d$, représenté par un $k$-vecteur $\bw$, le \emph{taux de dilatation} $\tau(W)$ de $W$ par le flot $(a_t)$ est le logarithme de la plus grande valeur propre apparaissant dans la décomposition de $\bw$ selon les espaces propres de $a_1$.
\end{defi}

\begin{rema}
Le taux de dilatation de $W$ s'obtient aussi simplement comme une limite:
\[
\tau(W) = \lim_{t\,\to\,\infty}\frac{1}{t}\log\left\|a_t\bw\right\|.
\]
Par conséquent, si $\Lambda=\Z v_1\oplus\dots\oplus\Z v_k$ est un réseau dans $W$, alors
\[
\covol_{a_tW}\left(a_t\Lambda\right) =_{t\,\to\,\infty} e^{t\left(\tau(W)+o(1)\right)}.
\]
\end{rema}

\begin{proof}[Démonstration du théorème~\ref{drapeau}]
Pour $i=1,\dots,d$, posons 
\[
\Lambda_i=\liminf \frac{1}{t}\log\lambda_i\left(a_t\Delta\right).
\]
On définit les indices $i_0,i_1,\,\dots,\,i_s,i_{s+1}$ par $i_0=0$, $i_{s+1}=d$ et
\[
\Lambda_1=\dots=\Lambda_{i_1}<\Lambda_{i_1+1}=\dots=\Lambda_{i_2}<\dots<\Lambda_{i_s+1}=\dots=\Lambda_d.
\]
Nous allons montrer par récurrence sur $k\in\se{0}{s}$ la propriété suivante: il existe un drapeau partiel $0=T_0<T_1<\dots<T_k$ de sous-espaces $\Delta$-rationnels tel que pour chaque $\ell\in\se{1}{k}$,
\begin{itemize}
\item $\dim T_\ell = i_\ell$;
\item pour tout $t>0$ assez grand, $a_tT_\ell$ contient les $i_\ell$ premiers minima de $a_t\Delta$;
\item si $i_{\ell-1}<i\leq i_\ell$, alors $\lim_{t\,\to\,\infty}\frac{1}{t}\log\lambda_i(a_t\Delta)=\Lambda_{i_\ell} = \frac{\tau(T_\ell)-\tau(T_{\ell-1})}{i_\ell-i_{\ell-1}}$.
\end{itemize}
De plus, pour tout $i\leq i_k$, $\lim_{t\to\infty}\frac{1}{t}\log\lambda_i(a_t\Delta)=\Lambda_i$.

Pour $k=0$, la propriété est triviale. Supposons donc la propriété vraie pour $k-1$, avec $k\in\se{1}{s+1}$. Soit $(t_n)$ une suite qui tend vers l'infini telle que $\Lambda_{i_k}=\lim\frac{1}{t_n}\log\lambda_{i_k}(a_{t_n}\Delta)$. Posant $\eps=\frac{1}{3}(\Lambda_{i_k+1}-\Lambda_{i_k})$, on a, pour $n$ suffisamment grand, $\lambda_{i_k}(a_{t_n}\Delta)\leq e^{t_n(\Lambda_{i_k}+\eps)}
\leq e^{t_n(\Lambda_{i_k+1}-2\eps)}
\leq e^{-t_n\eps}\lambda_{i_k+1}(a_{t_n}\Delta)$. D'après le théorème du sous-espace fort, il existe un sous-espace $\Delta$-rationnel $T_k$ de dimension~$i_k$ tel que le long d'une sous-suite de $(t_n)$, les $i_k$ premiers minima de $a_{t_n}\Delta$ sont toujours atteints dans $a_{t_n}T_k$.

Par hypothèse de récurrence, on a, pour $\ell<k$, si $i_{\ell-1}<i\leq i_\ell$,
\[
\lim\frac{1}{t_n}\log\lambda_i\left(a_{t_n}\Delta\right)=\Lambda_{i_\ell},
\]
et par définition de $(t_n)$ et des indices $i_{k-1}$ et $i_k$,
\[
\lim\frac{1}{t_n}\log\lambda_{i_k}(a_{t_n}\Delta) = \Lambda_{i_k} =\dots= \Lambda_{i_{k-1}+1}=\lim\frac{1}{t_n}\log\lambda_{i_{k-1}+1}\left(a_{t_n}\Delta\right).
\]
Comme $a_{t_n}T_k$ contient les $i_k$ premiers minima, cela implique
\begin{align*}
\covol_{a_{t_n}T_k}\left(a_{t_n}\left(T_k\cap\Delta\right)\right) &\leq e^{t_n\sum\limits_{\ell=1}^k\left(i_\ell-i_{\ell-1}\right)\Lambda_{i_\ell}+o(t_n)}
\\
\intertext{d'où l'on tire}
\tau(T_k) &\leq \sum_{\ell=1}^k (i_\ell-i_{\ell-1})\Lambda_{i_\ell}.
\end{align*}
Mais d'après l'hypothèse de récurrence, $\tau(T_{k-1}) = \sum_{\ell=1}^{k-1} (i_\ell-i_{\ell-1})\Lambda_{i_\ell}$, et donc
\begin{equation}\label{ikmoins}
\frac{\tau(T_k)-\tau(T_{k-1})}{i_k-i_{k-1}} \leq \Lambda_{i_k}=\dots=\Lambda_{i_{k-1}+1}.
\end{equation}
Comme l'hypothèse de récurrence implique $T_{k-1}\leq T_k$, le second théorème de Minkowski appliqué dans $a_tT_k$, donne, pour tout $t>0$ assez grand,
\begin{align*}
\lambda_{i_{k-1}+1}\left(a_t\Delta\right) \dots \lambda_{i_k}(a_t\Delta) & \ll \left(\frac{\covol\left(a_{t}T_k\right)}{\covol\left(a_{t}T_{k-1}\right)}\right)\\
& \leq \exp\big[t\left(\tau(T_k)-\tau\left(T_{k-1}\right)\right)+o(t)\big].
\end{align*}
Avec~\eqref{ikmoins}, cela montre qu'on doit avoir, pour chaque $i\in\left[i_{k-1}+1,i_k\right]$,
\[
\Lambda_i=\lim\frac{1}{t}\log\lambda_i(a_t\Delta)=\frac{\tau(T_k)-\tau(T_{k-1})}{i_k-i_{k-1}}.
\]
Pour $\eps=\tfrac{1}{3}(\Lambda_{i_k+1}-\Lambda_{i_k})$, pour $t>0$ assez grand, le sous-réseau $a_t(T_k\cap\Delta)$ contient $i_k$ vecteurs linéairement indépendants de norme inférieure à $e^{t(\Lambda_{i_k}+\eps)}$, et comme $\Lambda_{i_k+1}>\Lambda_{i_k}+2\eps$, tout vecteur de $a_t\Delta$ hors de $a_tT_k$ est de norme supérieure à $e^{t(\Lambda_{i_k}+2\eps)}$. Cela montre que pour $t>0$ assez grand, les $i_k$ premiers minima de $a_t\Delta$ sont atteints dans $a_tT_k$.
\end{proof}

Nous pouvons maintenant démontrer le théorème~\ref{roth_quadrique} annoncé dans l'introduction.

\begin{proof}[Démonstration du théorème~\ref{roth_quadrique}]
Tout d'abord, pour tout $x$ dans $X(\R)$, le corollaire~\ref{dirichlet} montre que $\beta(x)\geq 1$. De plus, s'il existe un sous-espace projectif rationnel totalement isotrope $W$ contenant $x$, on peut simplement considérer les approximations de $x$ par des points rationnels quelconques de $W$: le principe des tiroirs de Dirichlet appliqué dans $W$ montre alors que $\beta(x)\geq 1+\frac{1}{\dim W}$. Par conséquent, on a toujours $\beta(x)\geq 1+\frac{1}{k_x}$, et il nous suffira donc de démontrer ici l'inégalité réciproque.

Supposons tout d'abord que le point $x\in X(\QQ)$ est dans $\ker F$. Comme $\ker F$ est un sous-espace projectif rationnel, on a alors, pour tout $v\in X(\Q)\setminus\ker F$, $d(x,v)\gg H(v)^{-1}$. Par conséquent, toutes les très bonnes approximations de $x$ sont dans $\ker F$, et $\beta(x)$ est égal à l'exposant diophantien de $x$ dans le plan projectif $\ker F$. La formule désirée est donc un corollaire classique du théorème du sous-espace de Schmidt, pour lequel nous renvoyons à Schmidt~\cite[Corollary~1C]{schmidt_da}.

Nous supposons donc maintenant $x\not\in\ker F$. Par la proposition~\ref{dani}, l'exposant de $x$ est donné par
\[
\beta(x) = \frac{1}{1-\gamma(x)},
\]
où $\gamma(x)$ est le taux de fuite du flot $g_t^x\Z^d$ dans l'espace des réseaux. D'après le théorème~\ref{drapeau}, il existe un sous-réseau $T_1$ de rang $i_1$ dans $\Z^d$ tel que pour $t>0$ suffisamment grand, les $i_1$ premiers minima de $g_t^x\Z^d$ sont atteints dans $g_t^xT_1$. De plus, pour tout $\eps>0$, pour tout $t>0$ suffisamment grand,
\begin{equation}\label{kfirst}
e^{t\left(\frac{\tau(T_1)}{i_1}-\eps\right)}\leq \lambda_1\left(g_t^x\Z^d\right)\leq\dots\leq \lambda_{i_1}\left(g_t^x\Z^d\right) \leq e^{t\left(\frac{\tau(T_1)}{i_1}+\eps\right)}.
\end{equation}
Comme le flot diagonal considéré est $a_t=\diag(e^{-t},1,\,\dots,\,1,e^t)$, le taux de dilatation $\tau(T_1)$ ne peut prendre que les valeurs $0$ et $-1$.

Si $\tau(T_1)=0$, alors $\gamma(x)=0$ et $\beta(x)=1\leq 1+\frac{1}{k_x}$. Notons que dans ce cas, vues les observations faites au début de cette démonstration, $x$ ne saurait appartenir à un sous-espace totalement isotrope rationnel, et donc $k_x=+\infty$.

Si $\tau(T_1)=-1$, alors $\gamma(x)=\frac{1}{i_1}$, et donc $\beta(x)=1+\frac{1}{i_1-1}$. Nous allons voir que $T_1$ est un sous-espace totalement isotrope. En effet, pour $t>0$ suffisamment grand, \eqref{kfirst} montre qu'il existe une base $(u_1,\,\dots,\,u_{i_1})$ de $T_1\cap\Z^d$ telle que pour chaque $i$, $\|g_t^xu_i\|\leq e^{-t(\frac{1}{i_1}-\eps)}$. Comme $F$ est rationnelle, $F(u_i,u_j)$ est un rationnel à dénominateur borné, et nécessairement, en prenant $t>0$ assez grand $F(u_i,u_j)=F(g_t^xu_i,g_t^xu_j)=0$, i.e. $T_1$ est totalement isotrope. Le sous-espace $T_1$ est contracté par $g_t^x$, donc $u_xT_1$ contient la droite $[e_1]$ engendrée par le premier vecteur de la base canonique; de manière équivalente, $x=u_x^{-1}[e_1]$ appartient au sous-espace projectif totalement isotrope $\PP(T_1)$. Par suite, $k_x\leq i_1-1$, d'où $\beta(x)=1+\frac{1}{i_1-1}\leq 1+\frac{1}{k_x}$.
\end{proof}



\section{Non divergence pour \texorpdfstring{$\GO_F$}{GOF}}\label{sec:nd}

Soit $F$ une forme quadratique rationnelle sur $\R^d$, et $G=\GO_F(\R)<\GL_d(\R)$ le groupe de Lie réel des transformations linéaires de $\R^d$ qui préservent la forme quadratique $\Q$. Nous noterons $\Gamma=\GO_F(\Z)$ le sous-groupe de $G$ constitué des matrices à coefficients entiers dans la base canonique de $\R^d$. Comme $\Gamma$ est le stabilisateur de $\Z^d$ pour l'action de $G$ sur l'espace des réseaux de $\R^d$, on peut identifier $G/\Gamma$ à une partie fermée de l'espace des réseaux de $\R^d$:
\[
G/\Gamma \simeq \left\{g\Z^d\ ;\ g\in G\right\}.
\]
En particulier, pour $\Delta$ dans $G/\Gamma$, les minima successifs de $\Delta$ sont définis par
\[
\lambda_i(\Delta) = \inf\left\{\lambda>0\ \middle|\ \rg(\Delta\cap B(0,\lambda))\geq i\right\}.
\]

Étant donné un paramètre $K_1$, un espace métrique $X$ sera dit \emph{$K_1$-Besicovitch} si pour toute famille de boules ouvertes $B(a,r_a)$, $a\in A$, il existe une sous-famille dénombrable $(B(a_n,r_n))_{n\,\in\,\N}$ de multiplicité au plus $K_1$ telle que $A\subset\bigcup_{n\,\in\,\N}B(a_n,r_n)$. Si $B=B(x,r)$ est une boule ouverte dans $X$ et $m\geq 1$, nous noterons $mB=B(x,mr)$. Une mesure borélienne $\nu$ sur $X$ sera dite \emph{$K_2$-doublante} si pour toute boule ouverte $B\subset X$, $\nu(2B)\leq K_2\nu(B)$. Enfin, pour $C,\alpha>0$, une application $f:X\to\R$ sera dite \emph{$(C,\alpha)$-régulière} pour $\nu$ sur un ouvert $U\subset X$ si pour toute boule ouverte $B\subset U$ et tout $\eps\in]0,1]$,
\[
\nu\left(\left\{x\in B\ \middle|\ |f(x)|\leq\eps\right\}\right)
\leq C\left(\frac{\eps}{\sup_{x\,\in\,B\,\cap\,\Supp\nu}|f(x)|}\right)^\alpha \nu(B).
\]

Le but de cette partie est de démontrer le théorème ci-dessous, qui est l'analogue pour les groupes orthogonaux de~\cite[Theorem~0.2]{kleinbock_anextension}. Dans ce théorème et dans la suite de l'article, nous identifierons souvent un sous-réseau $\Lambda=\Z v_1\oplus\dots\oplus\Z v_k$ de $\Z^d$ au vecteur $\bw=v_1\wedge\dots\wedge v_k$ qui le représente dans $\wedge^k\R^d$, unique au signe près. En particulier, le covolume de $\Lambda$ dans l'espace vectoriel qu'il engendre s'obtient comme la norme $\|\bw\|$.

\begin{theo}[Non divergence pour $\GO_F$]\label{nondiv}
Soit $F$ une forme quadratique rationnelle sur $\R^d$, et $K_1,K_2,C_0,\alpha_0$ des paramètres strictement positifs. Il existe des constantes strictement positives $C$ et $\alpha$, fonctions de $K_1,K_2,C_0,\alpha_0$ et $F$ telles qu'on ait la propriété suivante.


On suppose que $X$ est un espace métrique $K_1$-Besicovitch, $\nu$ une mesure $K_2$-doublante sur $X$, $\rho\in]0,1]$, $h:X\to G$ une application continue et $B$ une boule ouverte telle que $5B\subset X$ et pour tout $k\in\{1,\,\dots,\,d\}$, pour tout sous-réseau de $\Z^d$ totalement isotrope pour $F$, représenté par un vecteur $\bw\in\wedge^k\Z^d$~:
\begin{enumerate}
\item $x\mapsto \|h(x)\bw\|$ est $(C_0,\alpha_0)$-régulière sur $5B$
\item \label{rhok}
$\sup_{x\,\in\,B\,\cap\,\Supp\,\nu}\|h(x)\bw\|\geq\rho^{k}$
\end{enumerate}
Alors, pour tout $\eps\in]0,1]$,
\[
\nu\left(\left\{x\in B\ \middle|\ \lambda_1\left(h(x)\Z^d\right)\leq\eps\rho\right\}\right)
\leq C\eps^\alpha\nu(B).
\]
\end{theo}

\begin{rema}
Si l'on se contentait d'appliquer le résultat de non divergence quantitative~\cite[Theorem~0.2]{kleinbock_anextension}, il serait nécessaire d'imposer la condition~\ref{rhok} à tous les sous-réseaux $\bw\leq\Z^d$. Le théorème ci-dessus montre qu'il suffit de considérer les sous-réseaux totalement isotropes pour $F$, ce qui nous sera utile dans la suite.
\end{rema}

\begin{rema}
Il est nécessaire d'imposer ici $\rho\in]0,1]$, car un supplémentaire rationnel $\bw$ de $\ker F$ vérifie $\|h(x)\bw\|\ll 1$ pour tout $x$ (noter que $\GO_{F'}(\R)$ préserve le volume si $F'$ est non dégénérée), et contient donc, par le premier théorème de Minkowski, un vecteur de norme bornée. Ceci est à comparer à l'énoncé analogue pour $\GL_d$, où l'on peut prendre $\rho$ arbitrairement grand~\cite[Theorem~2.2.1]{saxce_hnnondivergence}.
\end{rema}

Pour parvenir à ce résultat, notre démonstration s'éloigne sensiblement de la stratégie introduite par Margulis pour montrer ce type d'énoncé~\cite{dani_, kleinbockmargulis,margulis_}. Nous passerons par un résultat intermédiaire, le théorème~\ref{nondivk}, où l'hypothèse sur les fonctions $x\mapsto \|h(x)\bw\|$ ne porte que sur les sous-réseaux totalement isotropes de rang $k$ fixé. Nous aurons besoin du lemme suivant.

\begin{lemm}[Minima et espaces isotropes]\label{isotrope}
Soit $F$ une forme quadratique rationnelle sur $\R^d$, $G=\GO_F$ le groupe orthogonal associé, et $\Gamma=G(\Z)$. Il existe une constante $c>0$ telle que pour tout $\Delta$ dans $G/\Gamma$, si $(v_i)_{1\,\leq\,i\,\leq\,k}$ est une famille de vecteurs de $\Delta$ tels que pour chaque $i$, $\|v_i\|\leq c$, alors le sous-espace engendré par $(v_i)_{1\,\leq\,i\,\leq\,k}$ est totalement isotrope.
\end{lemm}

\begin{proof}
Soit $Q$ un dénominateur commun aux coefficients de $F$, et $\|F\|$ la norme de la forme bilinéaire $F$. Nous choisissons $c>0$ tel que $c^2\|F\|<Q^{-1}$. Il suffit de vérifier que $F(v_i,v_j)=0$ pour toute paire $(v_i,v_j)$ de vecteurs de la famille. D'une part, par notre choix de $c$, on a toujours $|F(v_i,v_j)|\leq \|F\|\|v_i\|\|v_j\| < Q^{-1}$. D'autre part, si on écrit $\Delta=g\Z^d$, avec $g\in G$, et $v_i=gu_i$, $u_i\in\Z^d$, alors, $F(v_i,v_j)=F(u_i,u_j)\in Q^{-1}\Z$. Donc $F(v_i,v_j)=0$.
\end{proof}



\begin{theo}[Non divergence pour les sous-réseaux de rang $k$]\label{nondivk}
Soit $F$ une forme quadratique rationnelle sur $\R^d$, $G=\GO_F$ le groupe orthogonal associé, $C_0,\alpha_0$ des paramètres strictement positifs, et $k\in\se{1}{d}$ \emph{fixé}. On suppose que $X$ est un espace métrique $K_1$-Besicovitch, $\nu$ une mesure $K_2$-doublante sur $X$, $\rho\in]0,1]$, $h:X\to G$ une application continue, et $B\subset X$ une boule ouverte telle que $5B\subset X$ et pour tout sous-réseau $\bw\leq\Z^d$ de rang $k$ totalement isotrope pour $F$:
\begin{enumerate}
\item $x\mapsto \|h(x)\bw\|$ est $(C_0,\alpha_0)$-régulière sur $5B$
\item \label{rhok'}
$\sup_{x\in B\cap\Supp\nu}\|h(x)\bw\|\geq\rho^k$
\end{enumerate}
Pour $\eps\in]0,1]$ et $M\geq 1$, on pose
\[
\cA_{M,\eps} = \left\{x\in B\ \middle|\
\begin{split}
\eps^M\rho &\leq\lambda_1\left(h(x)\Z^d\right)\\
\lambda_k\left(h(x)\Z^d\right)&\leq\eps\rho\\
\lambda_k\left(h(x)\Z^d\right)&\leq\eps\lambda_{k+1}\left(h(x)\Z^d\right)
\end{split}
\right\}.
\]
Alors,
\[
\nu(\cA_\eps)
\leq C_1M^k\eps^{\alpha_1}\nu(B),
\]
où $C_1=(2d)^kK_1K_2^3C_0$ et $\alpha_1=\frac{\alpha_0}{2}$.
\end{theo}

Le point crucial de la démonstration est le lemme suivant, dans lequel tous les minima jusqu'à $\lambda_k$ sont presque fixés.

\begin{lemm}\label{crucial}
Sous les hypothèses du théorème précédent, notons $c>0$ la constante donnée par le lemme~\ref{isotrope}. Soient $\rho_1,\dots,\rho_k\in]0,c[$ tels que $\rho_1\dots\rho_k\leq\eps^k\rho^k$. Pour $\eps\in]0,1]$, on pose
\[
\cA_\eps\left(\rho_1,\,\dots,\,\rho_k\right) = \left\{x\in B\ \middle|\
\begin{split}
\forall i\in\se{1}{k},\ \eps^{\frac{1}{2d}}\rho_i &\leq \lambda_i\left(h(x)\Z^d\right) \leq \rho_i\\
\lambda_k\left(h(x)\Z^d\right) &\leq\eps\lambda_{k+1}\left(h(x)\Z^d\right)
\end{split}
\right\}.
\]
Alors, pour tout $\eps>0$ suffisamment petit,
\[
\nu\left(\cA_\eps\left(\rho_1,\,\dots,\,\rho_k\right)\right)
\leq C_2\eps^{\alpha_2}\nu(B),
\]
où $C_2=K_1K_2^3C_0$ et $\alpha_2=\frac{\alpha_0}{2}$.
\end{lemm}

\begin{proof}
Pour plus de lisibilité, nous noterons $\lambda_i(x)=\lambda_i(h(x)\Z^d)$ et $S=\Supp\nu$. D'après le lemme~\ref{isotrope}, il existe pour chaque $x$ dans $\cA_\eps(\rho_1,\,\dots,\,\rho_k)$ un sous-réseau primitif totalement isotrope $\bw_x$ de rang $k$ tel que $\|h(x)\bw_x\|\leq\rho_1\dots\rho_k$. Soit $B_x$ une boule ouverte centrée en $x$ de rayon maximal tel que
\[
\forall y\in B_x\cap S,\ \left\|h(y)\bw_x\right\| \leq \eps^{-\frac{1}{2}}\rho_1\dots\rho_k.
\]
Noter que par continuité de $h$, la boule $B_x$ est de rayon strictement positif, et que $B_x\subset 3B$ et $2B_x\subset 5B$, car $\sup_{x\,\in\,B\,\cap\,S}\|h(y)\bw_x\|\geq\rho^k>\eps^{-\frac{1}{2}}\rho_1\dots\rho_k$. Soit $y\in S\cap B_x \cap\cA_\eps$ et $\bw$ un sous-réseau de rang $k$ non commensurable à $\bw_y$. Les minima successifs du sous-réseau $h(y)\bw$ sont minorés par $\lambda_1(y),\,\dots,\,\lambda_k(y)$, et comme $\bw$ n'est pas commensurable à $\bw_y$, le plus grand est minoré par $\lambda_{k+1}(y)$. Par le second théorème de Minkowski appliqué dans $\bw$, cela implique
\begin{align*}
\norm{h(y)\bw} &\gg \lambda_{k+1}(y)\lambda_{k-1}(y)\dots\lambda_1(y)
\\
\intertext{et donc}
\left\|h(y)\bw\right\| &\gg \frac{\lambda_{k+1}(y)}{\lambda_k(y)}\lambda_k(y)\dots\lambda_1(y)
\geq \eps^{-1}\eps^{\frac{k}{2d}}\rho_1\dots\rho_k > \|h(y)\bw_x\|,
\end{align*}
ce qui montre que $\bw_x=\bw_y$, et donc $\|h(y)\bw_x\|\leq\rho_1\dots\rho_k$. Comme l'application $y\mapsto\|h(y)\bw_x\|$ est $(C_0,\alpha_0)$-régulière sur $2B_x$ et vérifie $\sup_{y\in 2B_x\cap S}\|h(y)\bw_x\|\geq\eps^{-\frac{1}{2}}\rho_1\dots\rho_k$, on obtient
\begin{align*}
\nu\left(B_x\cap\cA_\eps\right) & \leq \nu\left(\left\{y\in B_x \ \middle|\ \left\|h(y)\bw_x\right\|\leq\rho_1\dots\rho_k\right\}\right)\\
& \leq C_0 \eps^{\frac{\alpha_0}{2}}\nu(2B_x)\\
& \leq C_0K_2 \eps^{\frac{\alpha_0}{2}}\nu(B_x).
\end{align*}
Par la propriété de Besicovitch, il existe une sous-famille dénombrable $(B_{x_i})_{i\,\in\,\N}$ de multiplicité au plus $K_1$ telle que $\cA_\eps\subset\bigcup_{i\,\in\,\N}B_{x_i}$, ce qui permet d'écrire
\[
\nu\left(\cA_\eps\right) \leq \sum_{i\,\in\,\N} \nu\left(B_{x_i}\cap\cA_\eps\right)
\leq C_0K_2\eps^{\frac{\alpha_0}{2}} \sum_{i\,\in\,\N} \nu\left(B_{x_i}\right)
\leq K_1K_2^3C_0\eps^{\frac{\alpha_0}{2}}\nu(B).
\]
\end{proof}

\begin{proof}[Démonstration du théorème~\ref{nondivk}]
Pour chaque famille d'entiers $m_i\in\{0,\dots,\linebreak2dM -1\}$, $i=1,\dots,k$, on pose
\begin{align*}
\cA_\eps^{\left(m_1,\,\dots,\,m_k\right)} &= \cA_\eps\left(\eps^{1+\frac{m_1}{2d}}\rho,\dots,\eps^{1+\frac{m_k}{2d}}\rho\right).
\\
\intertext{de sorte que}
\cA_{M,\eps} &\subset \bigcup_{m_1,\,\dots,\,m_k} \cA_\eps^{\left(m_1,\,\dots,\,m_k\right)}.
\end{align*}
D'après le lemme précédent, pour chaque $(m_i)_{1\,\leq\,i\,\leq\,k}$,
\begin{align*}
\nu\left(\cA_\eps^{\left(m_1,\,\dots,\,m_k\right)}\right)&\leq C_2\eps^{\alpha_2}\nu(B),
\\
\intertext{et donc}
\nu\left(\cA_{M,\eps}\right) &\leq (2dM)^kC_2\eps^{\alpha_2}\nu(B)
\end{align*}
ce qu'il fallait démontrer.
\end{proof}

Nous pouvons maintenant démontrer le théorème~\ref{nondiv}.

\begin{proof}[Démonstration du théorème~\ref{nondiv}]
Ici encore, on note $\lambda_1(x)=\lambda_1(h(x)\Z^d)$. Comme on peut toujours augmenter la constante $C$, on peut se contenter de démontrer l'inégalité du théorème pour $\eps>0$ suffisamment petit.

Il suffit de montrer que sous les hypothèses du théorème,
\begin{equation}\label{eps2eps}
\nu\left(\left\{x\in B\ \middle|\ \eps^2\rho \leq \lambda_1(x)\leq\eps\rho\right\}\right)
\leq C\eps^\alpha\nu(B).
\end{equation}
En effet, on aura alors, pourvu que $\eps<\frac{1}{2}$,
\begin{align*}
\nu\left(\left\{x\in B\ \middle|\ \lambda_1(x)\leq\eps\rho\right\}\right) & \leq \sum_{m\,\geq\,0} \nu\left(\left\{x\in B\ \middle|\ \eps^{2^{m+1}}\rho \leq \lambda_1(x)\leq\eps^{2^m}\rho\right\}\right)\\
& \leq C \sum_{m\,\geq\,0} \eps^{2^m\alpha}\nu(B)\\
& \leq \frac{C}{1-2^{-\alpha}} \eps^\alpha \nu(B).
\end{align*}
Pour voir~\eqref{eps2eps}, nous allons nous ramener au théorème~\ref{nondivk} avec pour paramètres $M=2$ et $\eps^{\frac{1}{2d}}$. Posons $s=\dim\ker F$. En appliquant les hypothèses du théorème au sous-réseau $\bw=\Z^d\cap\ker F$, avec l'égalité $\|h(x)\bw\|=\abs{\det h(x)}$, on trouve
\begin{equation}\label{det}
\nu\left(\left\{x\in B\ \middle|\ \abs{\det h(x)}\leq\eps^{\frac{1}{2d}}\rho^s\right\}\right) \leq C_0\eps^{\frac{\alpha_0}{2d}}\nu(B).
\end{equation}

Soit $x$ tel que $\lambda_1(x)\leq\eps\rho$ et $\abs{\det h(x)}\geq\eps^{\frac{1}{2d}}\rho^s$. On choisit $k\in\se{1}{d-1}$ minimal tel que $\lambda_k(x)\leq\eps^{\frac{1}{2d}}\lambda_{k+1}(x)$. Un tel $k$ existe nécessairement lorsque $\eps>0$ est suffisamment petit, sans quoi, partant de $\lambda_1(x)\leq\eps\rho$, on obtient, pour tout $k\geq 1$, $\lambda_k(x)\leq\eps^{1-\frac{k-1}{2d}}\rho$ puis $\lambda_1(x)\dots\lambda_d(x)\leq\eps\rho^d$, ce qui contredit $\lambda_1(x)\dots\lambda_d(x)\asymp \abs{\det h(x)}\geq \eps^{\frac{1}{2d}}\rho^s$. Pour ce choix de $k$, on doit avoir $\lambda_k(x)\leq\eps^{-\frac{k}{2d}}\lambda_1(x)\leq\eps^{\frac{1}{2}}\rho$.

Ainsi, l'ensemble des points $x$ tels que $\eps^2\rho\leq\lambda_1(x)\leq\eps\rho$ et $\abs{\det h(x)}\geq\eps^{\frac{1}{2d}}\rho^s$ est inclus dans la réunion
\[
\bigcup_{k=1}^{d-1}
\left\{x\in B\ \middle|\
\begin{split}
\eps^2\rho &\leq\lambda_1(x)\\
\lambda_k(x) &\leq\eps^{\frac{1}{2}}\rho\\
\lambda_k(x) &\leq\eps^{\frac{1}{2d}}\lambda_{k+1}(x)
\end{split}
\right\}.
\]
Pour $k$ fixé, le théorème~\ref{nondivk} appliqué avec $M=4d$ et $\eps^{\frac{1}{2d}}$ nous permet de majorer la mesure du $k^{\rm ième}$ ensemble de la réunion ci-dessus par
\[
C_1(4d)^k\eps^{\frac{\alpha_1}{2d}}\nu(B)
\]
puis, sommant sur $k$,
\[
\nu\left(\left\{x\in B\ \middle|\ \eps^2\rho\leq\lambda_1(x)\leq\eps\rho\ \text{et}\ \abs{\det h(x)}\geq\eps^{\frac{1}{2d}}\rho^s\right\}\right)
\leq C_1(4d)^d\eps^{\frac{\alpha_1}{2d}}\nu(B).
\]
Avec~\eqref{det}, cela donne bien, pour des constantes $C$ et $\alpha$ fonctions de $C_1$, $\alpha_1$ et $d$,
\[
\nu\left(\left\{x\in B\ \middle|\ \eps^2\rho\leq\lambda_1(x)\leq\eps\rho\right\}\right)
\leq C\eps^{\alpha}\nu(B).
\]
\end{proof}



\section{L'exposant diophantien d'une variété analytique}\label{sec:existence_exposant}


Dans toute cette partie $X$ désignera une quadrique projective rationnelle dans $\PP^{d-1}$, définie par une forme quadratique $F$. Nous voulons démontrer le théorème~\ref{exexpi} annoncé dans l'introduction, selon lequel l'exposant diophantien est constant presque partout sur une sous-variété analytique de $X(\R)$. L'argument est le même que dans l'article~\cite{kleinbock_dichotomy} de Kleinbock, bien que le cadre soit légèrement différent.

%\bigskip

Étant donnée une sous-variété analytique $\cM$ de dimension $m$ dans $X(\R)$, nous appellerons \emph{mesure de Lebesgue sur $\cM$} la restriction à $\cM$ de la mesure de Hausdorff de dimension $m$. Cette mesure sera la seule considérée sur $\cM$, et lorsque nous écrirons \og~pour presque tout $x$ dans $\cM$~\fg{}, ce sera toujours en référence à la mesure de Lebesgue sur $\cM$.


\begin{theo}[Exposant d'une mesure variété analytique]\label{exexp}
Soit $X$ une quadrique projective rationnelle dans $\PP^{d-1}$ telle que $X(\Q)\setminus\ker F\neq\varnothing$, et $\cM$ une sous-variété analytique connexe de $X(\R)$. Pour presque tout $x$ dans $\cM$,
\[
\beta(x) = \inf_{y\,\in\,\cM} \beta(y) = \bbeta(\cM).
\]
\end{theo}

La démonstration de ce théorème est une application de la proposition~\ref{dani} et du théorème~\ref{nondiv}. Pour chaque $x$ dans $X(\R)\setminus\ker F$ la proposition~\ref{dani} relie l'exposant diophantien $\beta(x)$ au taux de fuite d'une orbite diagonale $(a_tu_x\Z^d)$ dans l'espace des réseaux, où $u_x$ est un élément bien choisi de $\GO_F(\R)$. Nous voulons en premier lieu préciser le choix de $u_x$, de sorte que l'application
\begin{align*}
h:  X(\R)\setminus\ker F & \to  \GO_F(\R)\\
 x & \mapsto  u_x
\end{align*}
satisfasse localement les hypothèses du théorème~\ref{nondiv}.

Rappelons que le groupe $G=\GO_F(\R)$ agit naturellement à droite sur $X(\R)$ par la formule $[v]\cdot g = [g^{-1}v]$. D'après le théorème de Witt, cette action est transitive sur $X(\R)\setminus\ker F$. Comme $X(\Q)\setminus\ker F\neq\varnothing$, on peut choisir une base rationnelle $(e_i)_{1\,\leq\,i\,\leq\,d}$ de $\R^d$ dans laquelle la forme quadratique $F$ s'exprime
\[
F(x_1,\,\dots,\,x_d) = 2x_1x_d + G(x_2,\,\dots,\,x_{d-1}).
\]
Notons $P=\Stab[e_1]$ le sous-groupe qui stabilise la droite isotrope $[e_1]$, de sorte que $X(\R)\setminus\ker F$ s'identifie naturellement au quotient $P\bcs G$ via l'application
\begin{align*}
P\bcs G & \to  X(\R)\\
{[g]} & \mapsto  {\left[g^{-1}e_1\right]}
\end{align*}


\begin{rema}
Si $F$ est dégénérée, la variété $X\setminus\ker F$ n'est pas projective, mais seulement quasi-projective, car le sous-groupe $P$ n'est pas un sous-groupe parabolique de $G$. En effet, dans une décomposition $\R^d=\ker F\oplus V$, le groupe $G$ peut se voir sous forme matricielle par blocs comme $(
\begin{smallmatrix}\GL & *\\ 0 & \GO_{F'}
\end{smallmatrix})$.
Et l'on voit bien alors que $P$ ne contient pas le radical unipotent de $G$, constitué des matrices $(
\begin{smallmatrix} 1 & *\\ 0 & 1
\end{smallmatrix})$.
\end{rema}

Soit $U$ le sous-groupe unipotent opposé à $P$ ; ce sous-groupe est bien défini, car la projection de $P$ modulo le radical unipotent est un sous-groupe parabolique de la partie réductive de $G$. L'application de $U$ dans $X(\R)$ donnée par
\begin{align*}
\psi:  U & \to  X(\R)\\
 u & \mapsto  \left[u^{-1}e_1\right]
\end{align*}
est un difféomorphisme local au voisinage de $1$. Étant donné $x_0$ dans $X(\R)$, nous fixerons $g_0$ dans $G$ tel que $x_0=[g_0^{-1}e_1]$ pour définir, dans un voisinage de $x_0$,
\[
u_x = \psi^{-1}(x)g_0.
\]
Ensuite comme précédemment, pour $t>0$, on pose $a_t=\diag(e^{-t},1,\,\dots,\,1,e^t)$ et
\[
g_t^x = a_tu_x.
\]

\begin{rema}
En termes de l'action à droite de $G$ sur $X(\R)\simeq P\bcs G$, $u_x$ est l'unique élément de $Ug_{x_0}$ tel que $x=[e_1]\cdot u_x$. On remarquera que cette application n'est autre que la projection stéréographique, ou la carte affine usuelle, de $X(\R)$ au voisinage de $x_0$. En fait, pour les besoins de cet article, n'importe quelle section analytique locale $X(\R)=P\bcs G\to G$ au voisinage de $x_0$ conviendrait.
\end{rema}


\begin{prop}\label{regreg}
Soit $\cM$ une sous-variété analytique de $X(\R)$ et $\nu$ la mesure de Lebesgue sur $\cM$. Étant donné un point $x_0$ de $\cM$, il existe un voisinage $U_0$ de $x_0$ et des constantes $C_0,\alpha_0>0$ telles que pour tout $t>0$ et tout vecteur $\bw\in\wedge^k\R^d$, l'application
\[
x \mapsto \left\|g_t^x \bw\right\|
\]
est $(C_0,\alpha_0)$-régulière sur $U_0$ pour la mesure $\nu$.
\end{prop}

\begin{proof}
Soit $V$ un voisinage de $0$ dans $\R^m$ et $f:V\to\cM$ un paramétrage local de $\cM$ au voisinage de $x_0$, avec $f(0)=x_0$. Au voisinage de $x_0$, la mesure $\nu$ est équivalente à l'image de la mesure de Lebesgue sur $V$ par $f$, avec une densité continue et strictement positive. Nous voulons donc montrer que pour certaines constantes $C_0$ et $\alpha_0$, les applications $z\mapsto \|g_t^{f(z)}\bw\|^2$ sont $(C_0,\alpha_0)$-régulières pour la mesure de Lebesgue sur $V$.

Or, notant $(\be_I)_{|I|=k}$ les vecteurs de la base canonique de $\wedge^k\R^d$, $e^{\lambda_I t}$ la valeur propre de $a_t$ associée, et $\bw=\sum_I w_I\be_I$,
\begin{align*}
g_t^{f(z)}\bw &= \sum_I e^{\lambda_I t} w_I u_{f(z)}\be_I
\\
\intertext{et donc}
\left\|g_t^{f(z)}\bw\right\|^2 &= \sum_{I,J} e^{2\lambda_I t}w_I^2 \left\langle\be_J,u_{f(z)}\be_I\right\rangle^2.
\end{align*}
Cela montre que $z\mapsto\|g_t^{f(z)}\bw\|^2$ appartient au sous-espace vectoriel de dimension finie $\cF$ engendré par les applications $z\mapsto\bracket{\be_J,u_{f(z)}\be_I}^2$, avec $I,J\subset\se{1}{d}$. Comme toutes ces applications sont analytiques, on peut appliquer~\cite[Proposition~2.1]{kleinbock_dichotomy} et conclure qu'il existe des constantes $C_0,\alpha_0>0$ et un voisinage de $0$ sur lequel tout élément de $\cF$ est $(C_0,\alpha_0)$-régulier. C'est ce qu'on voulait.
\end{proof}


Nous pouvons maintenant démontrer le théorème~\ref{exexp}.

\begin{proof}[Démonstration du théorème~\ref{exexp}]
Notons $\nu$ la mesure de Lebesgue sur $\cM$, et fixons $x_0$ dans $\cM$. D'après la proposition~\ref{regreg}, il existe une boule ouverte $B$ centrée en $x_0$ et des constantes $C_0,\alpha_0$ telles que pour tout $t>0$ et tout $\bw$ dans $\wedge^k\R^d$, l'application $x\mapsto\|g_t^x\bw\|$ soit $(C_0,\alpha_0)$-régulière sur $5B$ par rapport à $\nu$.

Choisissons $y$ quelconque dans $B\cap\cM$, posons $\gamma=\gamma(y)$ et fixons $\eta>0$. Pour $t>0$ assez grand, $\lambda_1(g_t^y\Z^d)\geq e^{-(\gamma+\eta)t}$. Par le premier théorème de Minkowski, si $\bw\in\wedge^k\Z^d$ représente un sous-espace rationnel, on doit aussi avoir
\begin{align*}
\left\|g_t^y\bw\right\| \geq e^{-k(\gamma+\eta) t},
\\
\intertext{ce qui implique en particulier}
\sup_{x\,\in\,B\,\cap\,\cM} \left\|g_t^x\bw\right\| &\geq e^{-k(\gamma+\eta) t}.
\end{align*}
Le théorème~\ref{nondiv}, appliqué à $h:x\mapsto g_t^x$ avec $\rho=e^{-(\gamma+\eta)t}$ et $\eps=e^{-\eta t}$ montre que
\[
\nu\left(\left\{x\in B \ \middle|\ \lambda_1\left(g_t^x\Z^d\right)\leq e^{-\left(\gamma+2\eta\right)t}\right\}\right) \leq C e^{-\alpha\eta t}.
\]
Comme $\sum_{t\,\in\,\N}e^{-\alpha\eta t}<\infty$, le lemme de Borel--Cantelli implique que pour $\nu$-presque tout $x$ dans $B$, pour $t>0$ assez grand, $\lambda_1(g_t^x\Z^d)\geq e^{-(\gamma+2\eta)t}$, et donc $\gamma(x)\leq\gamma+2\eta$. En faisant tendre $\eta$ vers zéro, on obtient, pour presque tout $x$ dans $B$, $\gamma(x)\leq\gamma=\gamma(y)$, et comme $y$ est quelconque dans $B$, $\gamma(x)\leq\inf_{y\,\in\,B}\gamma(y)$. L'inégalité réciproque est évidente pour tout $x$ dans $B$, et par conséquent, pour $\nu$-presque tout $x$ dans $B$,
\[
\gamma(x) = \inf_{y\,\in\,B\,\cap\,\cM}\gamma(y).
\]
Ainsi, pour chaque $x_0$ dans $\cM$, il existe une boule $B$ sur laquelle l'exposant diophantien est constant presque partout, égal à $\ggamma(x_0)=\inf_{y\,\in\,B\,\cap\,\cM}\gamma(y)$. Comme l'application $x_0\mapsto\ggamma(x_0)$ est localement constante sur la variété connexe $\cM$, elle doit être constante, nécessairement égale à $\inf_{y\,\in\,\cM}\gamma(y)$. Avec la proposition~\ref{dani}, cela montre ce qu'on voulait.
\end{proof}

Cet énoncé, mis en lien avec le théorème~\ref{roth_quadrique}, donne la formule souhaitée pour l'exposant presque sûr d'un ensemble algébrique irréductible $\cM$ défini sur $\QQ$. Rappelons qu'un \emph{ensemble algébrique} $\cM$ est le lieu des zéros d'une famille de polynômes homogènes sur $X(\R)$. Il est dit \emph{irréductible} s'il n'est pas réunion de deux sous-ensembles algébriques stricts. Si $K$ est un sous-corps de $\C$, nous dirons aussi qu'un ensemble algébrique est \emph{défini sur $K$} si c'est le lieu des zéros d'une famille de polynômes à coefficients dans $K$. Comme dans le cas d'une variété analytique, nous munirons un ensemble algébrique $\cM$ de dimension $m$ de la restriction de la mesure de Hausdorff de dimension $m$.

\begin{rema}
Le support de la mesure de Lebesgue $\nu$ sur un ensemble algébrique $\cM$ n'est pas nécessairement dense dans $\cM$ pour la topologie usuelle. Des exemples sont donnés dans~\cite[Example~3.1.2]{bcr}. Cependant, tout point régulier (non singulier) de $\cM$ est inclus dans le support de $\nu$. Pour plus de détails sur les ensembles algébriques et leurs propriétés, nous renvoyons le lecteur à~\cite[Chapter~3]{bcr}.
\end{rema}

\begin{theo}[Exposant d'un ensemble algébrique défini sur $\QQ$]
Soit $X$ une quadrique rationnelle dans $\PP^{d-1}$. Si $\cM\subset X(\R)$ est un ensemble algébrique irréductible défini sur $\QQ$, alors, pour presque tout $x$ dans $\cM$,
\[
\beta(x) = 1 + \frac{1}{k_{\cM}},
\]
où $k_{\cM}$ est la dimension minimale d'un sous-espace projectif rationnel totalement isotrope contenant $\cM$.
\end{theo}


\begin{proof}
Comme dans la démonstration du théorème~\ref{roth_quadrique}, on voit facilement, par le principe des tiroirs de Dirichlet, que pour tout $x$ dans $\cM$, $\beta(x)\geq 1+\frac{1}{k_{\cM}}$. Cette égalité est bien sûr aussi valable lorsque $\cM$ n'est incluse dans aucun sous-espace rationnel isotrope, i.e. lorsque $k_{\cM}=+\infty$, puisqu'on a toujours $\beta(x)\geq 1$, d'après le corollaire~\ref{dirichlet}.

Soit maintenant $x_0$ un point régulier de $\cM$. Au voisinage de $x_0$, l'ensemble $\cM$ est une variété analytique. Pour voir que l'égalité a lieu presque sûrement dans un voisinage de $x_0$, il suffit donc, d'après le théorème~\ref{exexp}, de trouver des points $y$ arbitrairement proches de $x_0$ pour lesquels $\beta(y)=1+\frac{1}{k_{\cM}}$. Mais d'après le théorème~\ref{roth_quadrique}, cette égalité a lieu pour tout point algébrique $y$ tel que le plus petit sous-espace rationnel totalement isotrope contenant $y$ est égal au plus petit sous-espace rationnel totalement isotrope contenant $\cM$; et ces points sont denses dans $\cM$ au voisinage de $x_0$. Ainsi, l'égalité $\beta(x)=1+\frac{1}{k_{\cM}}$ est valable presque sûrement dans un voisinage de tout point régulier de $\cM$. Comme les points singuliers forment un sous-ensemble algébrique de codimension strictement positive dans $\cM$, les points réguliers sont de mesure pleine, et ce qui précède prouve le théorème.
\end{proof}



\section{Une formule pour l'exposant diophantien}\label{sec:formule_exposant}

Nous démontrons maintenant la formule pour l'exposant diophantien $\bbeta(\cM)$ d'un point pris aléatoirement sur une sous-variété analytique $\cM$ de $X(\R)$. Cette formule implique notamment que $\bbeta(\cM)$ ne dépend que du plus petit sous-espace totalement isotrope contenant $\cM$.

%\bigskip

Commençons par définir une hauteur sur les sous-espaces projectifs rationnels de $\PP^{d-1}$. Pour $k\in\se{1}{d}$, nous noterons $\Grass(k,d)$ la grassmannienne des sous-espaces vectoriels de dimension $k$ dans $\R^d$. Nous identifierons souvent dans la suite un sous-espace vectoriel de dimension $k$ dans $\R^d$ avec le sous-espace projectif de dimension $k-1$ correspondant dans $\PP^{d-1}$. Le \emph{plongement de Pl\"ucker} $p_k:\Grass(k,d)\hookrightarrow\PP(\wedge^k\R^d)$ est l'application qui associe à un sous-espace $W$ la droite $[\bw]$ dans $\wedge^k\R^d$, où $\bw$ est un vecteur représentant $W$, i.e. s'écrivant $\bw=v_1\wedge\dots\wedge v_k$ pour une certaine base $(v_1,\dots,v_k)$ de $W$. Le plongement de Pl\"ucker envoie les sous-espaces rationnels sur des droites rationnelles dans $\PP(\wedge^k\R^d)$, et cela permet de définir la hauteur d'un sous-espace projectif rationnel $W$ de dimension $k$ par la formule
\[
H(W) = H\left(p_k(W)\right),
\]
où la hauteur sur $\PP(\wedge^k(\Q^d))$ est la hauteur usuelle: si $\bw$ s'écrit dans la base $(e_I)_{\abs{I}=k}$ en coordonnées homogènes $\bw=[w_I]_{\abs{I}=k}$, avec les $w_I\in\Z$ premiers entre eux dans leur ensemble, alors $H(\bw)=\max_I\abs{w_I}$.

Soit maintenant $X$ une quadrique rationnelle dans $\PP^{d-1}$, et $S\subset X(\R)$. Si $W$ est un sous-espace projectif de $\PP^{d-1}(\R)$, nous noterons $d(S,W)=\sup_{x\,\in\,S} d(x,W)$. Rappelons que le $i$-ème exposant diophantien de $S$ est défini par
\[
\beta_i(S) = \inf\left\{\beta>0\ \middle|
\begin{split}
\exists c>0:\,\forall W\ \text{totalement isotrope rationnel}\\
\text{de dimension}\ i,\ d(S,W)\geq cH(W)^{-\beta}
\end{split}
\right\}.
\]
Le lecteur prendra garde au fait que dans cette définition $W$ est un sous-espace \emph{projectif}; sa dimension est donc inférieure d'une unité à celle du sous-espace vectoriel correspondant. Nous voulons en premier lieu démontrer que l'exposant d'une sous-variété analytique dans $X(\R)$ est toujours supérieur ou égal à la formule donnée par le théorème~\ref{formexpi}. C'est le contenu de la proposition suivante, dont la démonstration est une application du principe des tiroirs de Dirichlet.

\begin{prop}\label{dirichlet-fort}
Soit $X$ une quadrique rationnelle dans $\PP^{d-1}$, et $S$ une partie de $X(\R)$. Pour tout $x$ dans $S$,
\[
\beta(x) \geq \max\left(1,\max_{0\,\leq\,i\,\leq\,d-1} \frac{(i+1)\beta_i(S)}{1+i\beta_i(S)}\right).
\]
\end{prop}

\begin{proof}
D'après le corollaire~\ref{dirichlet}, on a toujours $\beta(x)\geq 1$. Il s'agit donc seulement de vérifier que si $x\in S$, on a bien pour chaque $i$, $\beta(x)\geq \frac{(i+1)\beta_i(S)}{1+i\beta_i(S)}$. Soit $\beta<\beta_i(S)$. On peut trouver un sous-espace projectif rationnel isotrope $W$ de dimension~$i$ et de hauteur arbitrairement grande tel que $d(S,W)\leq H(W)^{-\beta}$. En particulier, $d(x,W)\leq H(W)^{-\beta}$, i.e. il existe $x_0\in W$ tel que $d(x,x_0)\leq H(W)^{-\beta}$. Pour tout $Q>1$, le convexe
\[
C_Q=\left\{\bv\in W\ \middle|\ \norm{\bv}\leq Q\ \text{et}\ d_{\R^d}\left(x_0,\bv\right)\leq \tau\right\}
\]
est de volume dans $W$ égal à $\abs{C_Q}\asymp Q\tau^{i}$. Choisissant $\tau\gg H(W)^{\frac{1}{i}}Q^{-\frac{1}{i}}$, on obtient $\abs{C_Q}\gg H(W)$, et le premier théorème de Minkowski appliqué dans le réseau $W\cap\Z^d$, de covolume $H(W)$, montre qu'il existe un vecteur $\bv\neq 0$ tel que $\norm{\bv}\leq Q$ et
\[
d_{\R^d}\left(x_0,\bv\right)\ll H(W)^{\frac{1}{i}}Q^{-\frac{1}{i}}.
\]
Le vecteur $v\in\PP(W)(\Q)$ correspondant à $\bv$ vérifie alors $H(v)\leq Q$ et $d(x_0,v)\ll H(W)^{\frac{1}{i}}Q^{-\frac{1}{i}}H(v)^{-1}$, et donc
\begin{align*}
d(x,v) & \leq d\left(x,x_0\right) + d\left(x_0,v\right)\\
& \ll H(W)^{-\beta} + Q^{-\frac{1}{i}}H(W)^{\frac{1}{i}}H(v)^{-1}.
\end{align*}
On choisit $Q$ tel que $Q^{1+\frac{1}{i}}= H(W)^{\beta+\frac{1}{i}}$ pour obtenir
\[
d(x,v) \ll H(v)^{-\frac{(i+1)\beta}{1+i\beta}}.
\]
Cela montre que $\beta(x)\geq\frac{(i+1)\beta}{1+i\beta}$, et en faisant tendre $\beta$ vers $\beta_i(S)$, on trouve l'inégalité souhaitée.
\end{proof}

C'est pour majorer l'exposant d'une sous-variété analytique que nous utiliserons le théorème~\ref{nondiv}, à travers le lemme suivant.

\begin{lemm}\label{bc}
Soit $\cM$ une sous-variété analytique de $X(\R)$ et $x_0$ un point de $\cM$. Choisissons une boule $B$ centrée en $x_0$, et des constantes $C_0,\alpha_0$ telles que pour tout $t>0$ et tout sous-réseau $\bw$ de $\Z^d$, l'application $x\mapsto\|g_t^x\bw\|$ est $(C_0,\alpha_0)$-régulière sur $5B$. Étant donné $\gamma>0$, supposons qu'il existe $c>0$ tel que pour tout sous-réseau totalement isotrope $\Lambda$ de $\Z^d$, représenté par $\bw\in\wedge^k\Z^d$, pour tout $t>0$,
\[
\sup_{x\,\in\,B\,\cap\,\cM}\left\|g_t^x\bw\right\| \geq c e^{-k \gamma t}.
\]
Alors, pour presque tout $x$ dans $B\cap\cM$,
\[
\gamma(x) \leq \gamma.
\]
\end{lemm}


\begin{proof}
On peut supposer sans perte de généralité que $c\in]0,1]$. Soit $\eta>0$. Pour $t>0$, on applique le théorème~\ref{nondiv} à l'application $x\mapsto g_t^x\Z^d$, avec $\rho=ce^{-\gamma t}$ et $\eps=c^{-1}e^{-t\eta}$:
\[
\nu\left(\left\{x\in B\ \middle|\ \lambda_1\left(g_t^x\Z^d\right)\leq e^{-\left(\gamma+\eta\right)t}\right\}\right)
\leq Cc^{-\alpha} e^{-t\eta\alpha}.
\]
Par suite, $\sum_{t\,\in\,\N} \nu(\{x\in B\ |\ \lambda_1(g_t^x\Z^d)\leq e^{-(\gamma+\eta)t}\})<\infty$, et donc, d'après le lemme de Borel-Cantelli, pour $\nu$-presque tout $x$, pour $t$ assez grand, $\lambda_1(g_t^x\Z^d)\geq e^{-(\gamma+\eta)t}$, d'où $\gamma(x)\leq\gamma+\eta$. Comme ceci est valable pour tout $\eta>0$, on trouve bien, pour $\nu$-presque tout $x$ dans $B$, $\gamma(x)\leq\gamma$.
\end{proof}

Pour plus de clarté nous formulons aussi dès maintenant une estimation sur la norme d'un vecteur $g_t^x\bw$ qui nous sera utile ci-dessous.

\begin{lemm}\label{est}
Soit $x_0$ un point de $X(\R)$. Il existe une boule $B$ centrée en $x_0$ et une constante $c_0>0$ telle que pour tout $x\in B$, pour tout $t>0$, et pour tout sous-espace projectif $W$, représenté par un vecteur $\bw\in\wedge^k\Z^d$,
\[
\left\|g_t^x\bw\right\| \geq c_0\left(e^{-t} + d(x,W)\right)\|\bw\|.
\]
\end{lemm}

\begin{proof}
Décomposons $u_x\bw= \bw_x^{(-1)} + \bw_x^{(0)} + \bw_x^{(1)}$ selon les espaces propres de $a_t$, de sorte que
\[
g_t^x\bw = e^{-t}\bw_x^{(-1)} + \bw_x^{(0)} + e^t \bw_x^{(1)}.
\]
On a alors
\begin{align*}
\left\|g_t^x\bw\right\| & \asymp e^{-t}\left\|\bw_x^{(-1)}\right\| + \left\|\bw_x^{(0)}\right\| + e^t \left\|\bw_x^{(1)}\right\|\\
& \gg e^{-t}\left\|u_x\bw\right\| + \left\|\bw_x^{(0)}\right\| + \left\|\bw_x^{(1)}\right\|,
\end{align*}
et comme dans un voisinage de $x_0$, $u_x$ est uniformément bi-lipschitzien,
\[
\left\|g_t^x\bw\right\| \geq c\cdot\left(e^{-t}\left\|\bw\right\| + \left\|\bw_x^{(0)}\right\|+\left\|\bw_x^{(1)}\right\|\right).
\]
Soit $L$ le sous-espace vectoriel de $\wedge^k\R^d$ engendré par les $k$-plans contenant $e_1$, i.e. $L=\R e_1\wedge(\wedge^{k-1}\R^d)$. On a l'inégalité
\[
\left\|\bw_x^{(0)}\right\|+\left\|\bw_x^{(1)}\right\| \geq d\left(u_x\bw,L\right).
\]
Écrivons maintenant $e_1=v_1+e$, avec $v_1\in u_xW$ et $e\in (u_xW)^\perp$, de sorte que $\|e\|=d(e_1,u_xW)\asymp d(x,W)$. On adjoint à $v_1$ une famille de vecteurs unitaires pour obtenir une base orthogonale $(v_1,\,\dots,\,v_k)$ de $u_xW$. Alors,
\begin{align*}
\frac{1}{\left\|u_x\bw\right\|}u_x\bw & = \frac{v_1}{\|v_1\|}\wedge v_2\wedge\dots\wedge v_k\\
& = \frac{e_1}{\norm{v_1}}\wedge v_2\wedge\dots\wedge v_k - \frac{e}{\|v_1\|}\wedge v_2\wedge\dots\wedge v_k,
\end{align*}
et donc, comme $e$ est orthogonal à $e_1$ et à tous les $v_i$,
\[
\frac{1}{\left\|u_x\bw\right\|}d\left(u_x\bw, L\right) = \frac{1}{\|v_1\|}\|e\wedge v_2\wedge\dots\wedge v_k\| = \frac{\|e\|}{\|v_1\|} \gg d(x,W).
\]
Comme $u_x$ est uniformément bi-lipschitzien au voisinage de $x_0$, cela implique $d(u_x\bw,\linebreak L) \geq c_0 d(x,W)\|\bw\|$ et donc
\[
\left\|g_t^x\bw\right\| \geq c_0\left(e^{-t} + d(x,W)\right)\|\bw\|.
\]
\end{proof}


Nous pouvons enfin démontrer la formule pour l'exposant annoncée dans l'introduction, et que nous rappelons ici.

\begin{theo}[Formule pour l'exposant et hérédité]\label{formexp}
Soit $X$ une quadrique rationnelle dans $\PP^{d-1}$, et $\cM$ une sous-variété analytique connexe de $X(\R)$. L'exposant $\bbeta(\cM)$ est donné par la formule
\[
\bbeta(\cM) = \max\left(1,\max_{0\,\leq\,i\,\leq\,d-1} \frac{(i+1)\beta_i(\cM)}{1+i\beta_i(\cM)}\right).
\]
En particulier, $\bbeta(\cM)$ est déterminé par le plus petit sous-espace totalement isotrope $W_{\cM}$ contenant $\cM$.
\end{theo}


\begin{proof}
Comme la proposition~\ref{dirichlet-fort} montre que pour tout $x$ dans $\cM$,
\[
\beta(x) \geq \max\left(1,\max_{0\,\leq\, i\,\leq\,d-1} \frac{(i+1)\beta_i(S)}{1+i\beta_i(S)}\right).
\]
nous avons seulement à montrer l'inégalité réciproque pour presque tout $x$ dans $\cM$.

Soit $x_0$ un point de $\cM$, soit $B$ une boule centrée en $x_0$ telle que les lemmes~\ref{bc} et~\ref{est} s'appliquent, et posons $S=\cM\cap B$. Étant donné
\begin{equation}\label{max}
\beta> \max\left(1,\max_{0\,\leq\,i\,\leq\,d-1}\frac{(i+1)\beta_i(S)}{1+i\beta_i(S)}\right)
\end{equation}
nous voulons montrer que pour presque tout $x$ dans un voisinage de $x_0$, $\beta(x)\leq\beta$. Par la proposition~\ref{dani}, il suffit de voir que pour presque tout $x$ dans un voisinage de $x_0$, $\gamma(x)\leq\gamma=1-\frac{1}{\beta}$, et d'après le lemme~\ref{bc}, il suffit de vérifier qu'il existe $c>0$ tel que pour tout sous-réseau totalement isotrope $\bw\in\wedge^k\Z^d$, pour tout $t>0$,
\[
\sup_{x\,\in\,S}\left\|g_t^x\bw\right\| \geq c e^{-k\gamma t}.
\]
Or, d'après le lemme~\ref{est},
\[
\sup_{x\,\in\,S}\left\|g_t^x\bw\right\| \gg e^{-t}\|\bw\| + d(S,\bw)\|\bw\|.
\]
Si $k\gamma\geq 1$, on a pour tout $t>0$, $\|g_t^x\bw\|\geq e^{-t}\geq e^{-k\gamma t}$. Sinon, $k\gamma< 1$; choisissant $i=k-1$ dans le maximum de l'inégalité~\eqref{max}, on trouve
\[
\frac{\beta}{k-(k -1)\beta}>\beta_{k-1}(S),
\]
donc, par définition de $\beta_{k-1}(S)$, pour un certain $c>0$ indépendant de $\bw$,
\[
d(S,\bw) \geq c \|\bw\|^{-\frac{\beta}{k-(k-1)\beta}} = c \|\bw\|^{\frac{-1}{1-k\gamma}}.
\]
Ceci implique
\begin{align*}
e^{-t}\|\bw\| + d(S,\bw)\|\bw\| & \gg e^{-t}\|\bw\|+\|\bw\|^{\frac{-k\gamma}{1-k\gamma}}\\
& \geq \left(e^{-t}\|\bw\|\right)^{k\gamma}\|\bw\|^{-k\gamma} = e^{-k\gamma t}.
\end{align*}
Ainsi, pour presque tout $x$ dans un voisinage de $x_0$ dans $\cM$,
\[
\beta(x) = \max\left(1,\max_{0\,\leq\,i\,\leq\,d-1}\frac{(i+1)\beta_i(S)}{1+i\beta_i(S)}\right).
\]
D'après le lemme~\ref{her} ci-dessous, les exposants $\beta_i(S)$ ne dépendent que du plus petit sous-espace totalement isotrope $W_S$ contenant $S$. Or, comme $\cM$ est analytique, toute partie $S\subset\cM$ contenant un ouvert de $\cM$ vérifie $W_S=W_{\cM}$, et donc pour chaque $i$, $\beta_i(S)=\beta_i(\cM)$.
\end{proof}

Il nous reste seulement à démontrer le lemme d'hérédité suivant. Rappelons qu'étant données deux parties $S$ et $T$ dans $\PP^{d-1}(\R)$, on note $d(S,T)=\sup_{x\,\in\,S}d(x,T)$.

\begin{lemm}\label{her}
Soit $X$ une quadrique rationnelle dans $\PP^{d-1}$, et $S\subset X(\R)$. Notons $W_S$ le plus petit sous-espace projectif totalement isotrope contenant $S$, ou $W_S=X(\R)$ si $S$ n'est inclus dans aucun sous-espace totalement isotrope. Il existe une constante $c>0$ telle que pour tout sous-espace projectif totalement isotrope $W$,
\[
c\cdot d(W_S,W) \leq d(S,W) \leq d(W_S,W).
\]
En particulier, pour chaque $i\in\se{0}{d-1}$, $\beta_i(S)=\beta_i(W_S)$.
\end{lemm}


\begin{proof}
L'inégalité de droite est claire, puisque $S$ est inclus dans $W_S$. Pour l'inégalité de gauche, on peut supposer $\dim W=i$. On distingue alors deux cas.\\
\noindent{\underline{Premier cas:}} $S$ n'est inclus dans aucun sous-espace totalement isotrope de dimension~$i$. Alors, par compacité de la grassmannienne des sous-espaces totalement isotropes de dimension~$i$, il existe $c>0$ tel que pour tout $W$ totalement isotrope de dimension~$i$, $d(S,W)\geq c$. Cela implique ce que l'on veut.

\noindent{\underline{Second cas:}} $\dim W_S\leq i$. Alors $W_S$ co{\"\i}ncide avec l'adhérence linéaire de $S$ dans $\PP^{d-1}(\R)$. Sur $\wedge^{i+1}\R^d$, on définit deux semi-normes $N_1$ et $N_2$ par
\[
N_1(\bw)=\sup_{x\,\in\,S} \|x\wedge\bw\| \quad\text{et}\quad N_2(\bw)=\sup_{x\,\in\,W_S}\|x\wedge\bw\|.
\]
(Dans les formules ci-dessus, on identifie $x\in\PP^{d-1}(\R)$ avec un représentant unitaire dans $\R^d$.) Ces deux semi-normes ont le même noyau, et sont donc équivalentes: il existe une constante $c>0$ telle que pour tout $\bw$ dans $\wedge^{i+1}\R^d$, $N_1(\bw)\geq cN_2(\bw)$. Cela montre ce qu'on voulait, car si $W$ est un sous-espace projectif de dimension~$i$, représenté par un vecteur unitaire $\bw$ dans $\wedge^{i+1}\R^d$, alors $N_1(\bw)=d(S,W)$ et $N_2(\bw)=d(W_S,W)$.
\end{proof}

Le corollaire~\ref{petitplus}, qui renforce l'énoncé du théorème~\ref{algi}, se déduit de ce résultat plus faible, à l'aide du théorème~\ref{formexpi}.

\begin{proof}[Démonstration du corollaire~\ref{petitplus}]
D'après le théorème~\ref{formexpi}, l'exposant presque sûr $\bbeta(\cM)$ est entièrement déterminé par $W_{\cM}$, et donc $\bbeta(\cM)=\bbeta(W_{\cM})$. Mais par ailleurs, le théorème~\ref{algi} appliqué à la variété $\cM'=W_{\cM}$ montre que $\bbeta(W_{\cM})=1+\frac{1}{k_{\cM}}$. C'est le résultat souhaité.
\end{proof}


\section{Conclusion}

Nous avons montré que certains théorèmes d'approximation diophantienne sur l'espace projectif admettent des analogues (en fait, se généralisent) sur les quadriques projectives rationnelles. Il est souvent facile d'imaginer la version pour les quadriques d'un théorème pour l'espace projectif, mais la démonstration peut requérir l'introduction de nouvelles méthodes. Par exemple, nous avons eu besoin ici d'un énoncé de non divergence pour le groupe orthogonal $\GO_F(\R)$, qui ne semble pas se déduire des résultats de Kleinbock et Margulis, et pour lequel nous avons dû trouver une autre approche.

La généralité du théorème~\ref{nondiv} permet de montrer que les théorèmes~\ref{exexpi}, \ref{algi}, \ref{formexpi} et le corollaire~\ref{petitplus} sont encore valables pour toute mesure $\nu$ qui satisfait la conclusion de la proposition~\ref{regreg}. Comme cela a été fait par Kleinbock, Lindenstrauss et Weiss~\cite{klw} dans le cadre de l'espace projectif, il serait naturel de chercher à obtenir un critère géométrique simple sur $\nu$ pour que cette condition soit vérifiée. Disons qu'une mesure borélienne $\nu$ sur la quadrique rationnelle $X(\R)$ est \emph{localement régulière} en un point $x_0$ s'il existe un voisinage $U_0$ de $x_0$ et des constantes strictement positives $C_0$ et $\alpha_0$ telles que pour tout sous-espace projectif $W$ totalement isotrope rationnel, pour toute boule ouverte $B\subset U_0$,
\[
\nu\left(\left\{x\in B\ \middle|\ d(x,W)\leq\eps\right\}\right) \leq C_0\left(\frac{\eps}{\sup\limits_{x\,\in\,B\,\cap\,\Supp\nu}d(x,W)}\right)^{\alpha_0}\nu(B).
\]
Il serait naturel d'espérer le résultat suivant.

\begin{conj}
Soit $X$ une quadrique rationnelle dans $\PP^{d-1}$, et $\nu$ une mesure localement régulière en un point $x_0\in X(\R)$. Alors, il existe un voisinage $U_0$ de $x_0$ tel que pour $\nu$-presque tout $x$ dans $U_0$,
\[
\beta(x) = \inf_{y\,\in\,U_0\,\cap\,\Supp\nu}\beta(y).
\]
\end{conj}

Malheureusement, nous n'avons pas su obtenir dans ce cadre un analogue satisfaisant de la proposition~\ref{regreg}. Le résultat que nous avons obtenu, qui met en jeu d'autres sous-variétés linéaires de la quadrique, ne nous a pas semblé assez naturel pour être inclus ici, bien qu'il s'applique à une classe de mesures sensiblement plus générale que celle considérée dans le théorème~\ref{exexp}.



\bibliography{Desaxce}
\end{document}
